\section{Homological algebra}\label{ha:sec}

\providecommand{\fA}{\mathfrak A}
\providecommand{\hot}{\mathbin{\widehat\otimes}}
\providecommand{\Sh}{\operatorname{Sh}}
\providecommand{\Om}{\Omega}
\providecommand{\dOm}{\dot\Omega}

This appendix proves the homological background facts (B4), (B5), (B6)
and (B9) of Section~\ref{subsec:background}, and it analyses the convergence of the
order spectral sequence, which Remark~\ref{rem:poles} only summarises. The
arguments are adapted from Wodzicki \cite{Wodzicki}, from
\cite[Ch.~10]{MelroseMicro} and from the companion notes \cite{HHb}; we
say in each place what is taken from where. The status of each item is as
follows.
\begin{itemize}
\item (B4), Wodzicki excision: proved in full
(Theorem~\ref{ha:thm:excision}), using three stated facts about nuclear
spaces, (F1)--(F3) below.
\item (B5), Connes' theorem: proved for \(\CI(M)\), \(M\) compact with
corners (Theorem~\ref{ha:thm:connes}), and for the formal, Laurent and
flat coefficient algebras of the notes (Theorem~\ref{ha:thm:variants}).
\item (B6): the Künneth formula is proved in the form used
(Proposition~\ref{ha:prop:kunneth}); the Eilenberg--Zilber theorem is
quoted. Morita invariance for matrices, smoothing-kernel algebras and
full corners is proved (Lemma~\ref{ha:lem:morita},
Proposition~\ref{ha:prop:corner}).
\item Convergence: the algebraic convergence proposition is proved
(Proposition~\ref{ha:prop:algconv}), and so is convergence for the
closed symbol algebra and for \(\Is\) (Theorem~\ref{ha:thm:conv}). For
\(\As\) and \(\Ads\) the b-calculus argument of \cite{HHb} does
\emph{not} carry over, because of the cusp twist
(Example~\ref{ha:ex:obstruction}). We prove what survives
(Proposition~\ref{ha:prop:laurent}), we show that convergence for \(\As\)
and convergence for \(\Ads\) are equivalent
(Proposition~\ref{ha:prop:equiv}), and we prove both for \(n=1\) by the
\(x\)-adic filtration (Proposition~\ref{ha:prop:xadic},
Corollary~\ref{ha:cor:n1}). For every \(n\) convergence is proved in
Appendix~\ref{su:sec} (Theorem~\ref{su:thm:B10}), following Lauter and
Moroianu \cite{LauterMoroianu}; this is item (B10). The \(x\)-adic
spectral sequence of \(\Id\) converges (Proposition~\ref{ha:prop:Idconv}).
All of this uses the mixed complexes of Lemma~\ref{ha:lem:mixed}, which is
proved in full.
\item The boundary connecting map of \cite[Prop.~5(3)]{MN}, item (B9), is
proved by a leading-symbol argument with a linear lift
(Lemma~\ref{ha:lem:connlead}, Theorem~\ref{ha:thm:bdry}).
\end{itemize}

\subsection{Chains and three facts about nuclear spaces}\label{ha:subsec:facts}

We use the chain model of Section~\ref{subsec:chains}. A \emph{filtered
nuclear algebra} is an increasing union \(\fA=\bigcup_{N\ge0}\fA_N\) of
nuclear Fr\'echet spaces, each closed in the next, such that
multiplication maps \(\fA_N\times\fA_M\to\fA_{N+M}\) continuously. Its
chain spaces are \(C_k(\fA)=\varinjlim_N\fA_N^{\hot(k+1)}\), and \(d_i\),
\(b\), \(b'\) are as in Section~\ref{subsec:chains}. A nuclear Fr\'echet
algebra is the case \(\fA_N=\fA\). A closed ideal \(J\) and the quotient
\(\fA/J\) have the pieces \(J_N=J\cap\fA_N\) and \(\fA_N/J_N\). Homology
is always taken algebraically. Filtered colimits of vector spaces are
exact, so homology commutes with \(\varinjlim_N\). In the proofs below we
therefore work with one level \(N\) at a time and suppress \(N\), except
where a map changes the level. We use the following facts
\cite[Chs.~43, 50, 51]{Treves}.
\begin{enumerate}
\item[(F1)] Closed subspaces and Hausdorff quotients of nuclear Fr\'echet
spaces are nuclear Fr\'echet. If \(E\) is nuclear Fr\'echet and
\(0\to F'\to F\to F''\to0\) is an exact sequence of Fr\'echet spaces, then
\(0\to E\hot F'\to E\hot F\to E\hot F''\to0\) is exact.
\item[(F2)] (Kernel theorem.) For compact manifolds with corners
\(\CI(M)\hot\CI(N)=\CI(M\times N)\). The same holds for \(\dCI\), and
\(\cS(\R^a)\hot\cS(\R^b)=\cS(\R^{a+b})\).
\item[(F3)] If \(E\) is nuclear Fr\'echet, then \(E\hot\prod_jF_j=
\prod_jE\hot F_j\) for countable products of Fr\'echet spaces. For
example \(\CI(N)[[x]]^{\hot(k+1)}=\CI(N^{k+1})[[x_0,\dots,x_k]]\).
\end{enumerate}

\begin{lemma}[Exactness of chain spaces]\label{ha:lem:exact}
Let \(J\subset\fA\) be a closed ideal. For \(0\le s\le n\) put
\[
 H^s_n=\sum_{i\le s}\fA^{\hot i}\hot J\hot\fA^{\hot(n-i)}
 \subset\fA^{\hot(n+1)},
\]
the algebraic sum of the images, which are injective by (F1). Put
\(H^s_n=H^n_n\) for \(s>n\). Then:
\begin{enumerate}
\item \(\fA^{\hot(n+1)}\to(\fA/J)^{\hot(n+1)}\) is onto, with kernel
\(H^n_n\);
\item \(H^s_n/H^{s-1}_n\cong(\fA/J)^{\hot s}\hot J\hot\fA^{\hot(n-s)}\) for
\(1\le s\le n\);
\item \(J^{\hot q}\hot\fA^{\hot r}/J^{\hot(q+1)}\hot\fA^{\hot(r-1)}\cong
J^{\hot q}\hot(\fA/J)\hot\fA^{\hot(r-1)}\) for \(q,r\ge1\).
\end{enumerate}
The same holds for filtered nuclear algebras, with \(\varinjlim_N\)
applied throughout.
\end{lemma}

\begin{proof}
(3) is (F1) applied to \(0\to J\to\fA\to\fA/J\to0\) tensored with nuclear
factors. For (1) and (2), let \(\pi_s:\fA^{\hot(n+1)}\to(\fA/J)^{\hot s}
\hot\fA^{\hot(n+1-s)}\) project the first \(s\) factors. We show by
induction that \(\ker\pi_s=H^{s-1}_n\). For \(s=1\) this is (F1). Write
\(\pi_{s+1}=\pi'\circ\pi_s\), where \(\pi'\) projects the factor in slot
\(s\). By (F1), \(\ker\pi'=(\fA/J)^{\hot s}\hot J\hot\fA^{\hot(n-s)}\),
and \(\pi_s\) maps \(\fA^{\hot s}\hot J\hot\fA^{\hot(n-s)}\) onto it.
Hence \(\ker\pi_{s+1}=\ker\pi_s+\fA^{\hot s}\hot J\hot\fA^{\hot(n-s)}=
H^s_n\). The case \(s=n+1\) is (1), and surjectivity is (F1) again. For
(2), \(\pi_s\) restricted to \(H^s_n\) has kernel \(H^{s-1}_n\) and image
\(\pi_s(\fA^{\hot s}\hot J\hot\fA^{\hot(n-s)})=(\fA/J)^{\hot s}\hot J\hot
\fA^{\hot(n-s)}\).
\end{proof}

In particular \(\cK_n=H^n_n\) in \eqref{eq:chainses}. No continuous linear
section of \(\fA\to\fA/J\) is needed; the Taylor maps of the notes have
none.

\begin{definition}\label{ha:def:Hunital}
\(J\) is \emph{topologically H-unital} if there are continuous maps
\(E:J_N^{\hot(k+1)}\to J_N^{\hot(k+2)}\), for all \(N\) and \(k\ge0\), with
\(b'E+Eb'=1\) (where \(b'=0\) on \(C_0\)). We call \(E\) an \emph{extra
degeneracy} if moreover \(d_0E=1\) and \(d_{i+1}E=Ed_i\) for \(0\le i<k\);
then \(b'E+Eb'=1\) holds automatically.
\end{definition}

A unital algebra has the extra degeneracy \(E(c)=1\otimes c\). The
contractions of Lemmas~\ref{lem:Hunital} and~\ref{lem:flatHunital} are
extra degeneracies. If \(J\) is topologically H-unital, then the bar
complex is acyclic in degree \(0\) as well: \(b'E=1\) on \(C_0(J)\) says
that multiplication \(J\hot J\to J\) is onto. If \(E_1\), \(E_2\) are extra
degeneracies of \(J_1\), \(J_2\), then \(E_1\otimes E_2\) is an extra
degeneracy of \(J_1\hot J_2\), because the face maps of \(J_1\hot J_2\)
are the tensor products of the face maps.

\subsection{Wodzicki's excision theorem (B4)}\label{ha:subsec:excision}

Throughout this subsection \(J\subset\fA\) is a closed ideal of a filtered
nuclear algebra, and \(J\) is topologically H-unital with contraction
\(E\). We start with the case in which excision fails.

\begin{example}[Failure without H-unitality]\label{ha:ex:noexcision}
Let \(\fA=\C[x]\) (with algebraic tensor products) and \(J=x\C[x]\). Here
\(J\) is not H-unital, since \(J^2=x^2\C[x]\ne J\). The algebras are
commutative, so \(\HH_1(\fA)=\C[x]\,dx\) by HKR
(Theorem~\ref{ha:thm:poly}) and \(\HH_1(\fA/J)=\HH_1(\C)=0\). If excision
held, \(\HH_1(J)\to\HH_1(\fA)\) would be onto. But the HKR form of a chain
\(\sum a\otimes c\) with \(a,c\in J\) is \(\sum a\,dc\in x\C[x]\,dx\), so
the image misses \([dx]\).
\end{example}

The proof has two steps. Let \(\cK'_n=J\hot\fA^{\hot n}\) be the chains
whose first entry lies in \(J\). It is a subcomplex of \(\cK\): the faces
\(d_0\) and \(d_n\) put \(a_0a_1\) or \(a_na_0\) in slot \(0\), and both
lie in \(J\). We show that \(C(J)\subset\cK'\) and \(\cK'\subset\cK\) are
quasi-isomorphisms. The first step uses the bar complex of \(J\) directly.
The second needs it with coefficients.

\begin{lemma}[Bar complex with coefficients]\label{ha:lem:barW}
Let \(W\) be a left \(\fA\)-module which is a filtered nuclear space, with
continuous action \(\fA_N\times W_M\to W_{N+M}\), such that \(J\cdot W=0\).
Put \(P_m=J\hot\fA^{\hot m}\hot W\) for \(m\ge0\), with
\[
 b'_W(a_0\otimes\dots\otimes a_m\otimes w)=\sum_{i=0}^{m-1}(-1)^i
 a_0\otimes\dots\otimes a_ia_{i+1}\otimes\dots\otimes w+(-1)^m
 a_0\otimes\dots\otimes a_{m-1}\otimes a_mw .
\]
Then \((P_\bullet,b'_W)\) is acyclic, in degree \(0\) included.
\end{lemma}

\begin{proof}
\(P\) is the two-sided bar complex of the right module \(J\) and the left
module \(W\), so \(b_W'^2=0\). For \(0\le r\le m\) let \(G^r_m=J^{\hot(m+1-r)}
\hot\fA^{\hot r}\hot W\), the chains whose first \(m+1-r\) entries lie in
\(J\), and put \(G^r_m=P_m\) for \(r>m\). Then \(G^r\) is a subcomplex. A
face inside the initial \(J\)-block keeps the block, of length \(m-r\) in
degree \(m-1\), so it lands in \(G^r_{m-1}\). The face joining the last
block entry to the next entry gives an element of \(J\), and lands in
\(G^{r-1}_{m-1}\). A face joining two later entries, or the module face,
keeps the block and lands in \(G^{r-1}_{m-1}\). If \(r=0\) the module face
is \(a_mw=0\), because \(a_m\in J\).

So \(G^0=(C_\bullet(J),b')\hot W\) is acyclic, with contraction
\(E\otimes1\). For \(r\ge1\), Lemma~\ref{ha:lem:exact}(3) gives
\(G^r_m/G^{r-1}_m\cong J^{\hot(m+1-r)}\hot(\fA/J)\hot\fA^{\hot(r-1)}\hot W\).
By the list above, only the faces inside the block survive in the
quotient. Hence \(G^r/G^{r-1}\) is the bar complex of \(J\) (in degrees
\(m-r\ge0\)) tensored with spectators. Let \(\tilde E\) apply \(E\) to the
block. It maps \(G^{r-1}\) to \(G^{r-1}\), since it lengthens the block,
and \(b\tilde E+\tilde Eb\equiv(b'E+Eb')\otimes1=1\) modulo \(G^{r-1}\).
So each \(G^r/G^{r-1}\) is acyclic. Since \(G^{m+1}_j=P_j\) for \(j\le
m+1\), we get \(H_m(P)=H_m(G^{m+1})\cong H_m(G^0)=0\).
\end{proof}

\begin{proposition}[First step]\label{ha:prop:stepA}
The inclusion \(C(J)\to\cK'\) is a quasi-isomorphism.
\end{proposition}

\begin{proof}
Filter \(\cK'\) as in the previous proof: \(G^r_n=J^{\hot(n+1-r)}\hot
\fA^{\hot r}\) for \(r\le n\), and \(G^r_n=\cK'_n\) for \(r>n\). Then
\(G^0=C(J)\). The faces behave as before. The cyclic face \(d_n\) puts
\(a_na_0\in J\) in front and keeps the block, so it lands in
\(G^{r-1}_{n-1}\) when \(r\ge1\), and in \(G^0_{n-1}\) when \(r=0\). The
quotients \(G^r/G^{r-1}\), \(r\ge1\), are bar complexes of \(J\) tensored
with spectators. They are acyclic by the contraction \(\tilde E\) of the
previous proof. Since \(G^r_n=\cK'_n\) for \(r\ge n\), we get
\(H_n(C(J))=H_n(G^0)\cong H_n(G^{n+1})=H_n(\cK')\).
\end{proof}

\begin{proposition}[Second step]\label{ha:prop:stepB}
The inclusion \(\cK'\to\cK\) is a quasi-isomorphism.
\end{proposition}

\begin{proof}
Use the subspaces \(H^s_n\) of Lemma~\ref{ha:lem:exact}: chains with an
entry from \(J\) in one of the first \(s+1\) slots. Then \(H^0=\cK'\) and
\(H^n_n=\cK_n\). Each \(H^s\) is a subcomplex. Indeed, let \(a_i\in J\)
with \(i\le s\). A face \(d_j\) with \(j<i\) moves it to slot \(i-1\); the
faces \(d_{i-1}\), \(d_i\) and \(d_n\) produce an element of \(J\) in a
slot \(\le i\); the other faces leave slot \(i\) alone. Fix \(s\ge1\). By
Lemma~\ref{ha:lem:exact}(2),
\[
 H^s_n/H^{s-1}_n\cong(\fA/J)^{\hot s}\hot J\hot\fA^{\hot(n-s)} .
\]
Write an element as \(\bar a_0\otimes\dots\otimes\bar a_{s-1}\otimes a_s
\otimes\dots\otimes a_n\) with \(a_s\in J\). The faces \(d_j\), \(j\le
s-1\), produce an element of \(J\) in a slot \(\le s-1\), so they vanish
on the quotient. The faces \(d_j\), \(s\le j\le n-1\), multiply
neighbours in the segment \((a_s,\dots,a_n)\). The cyclic face puts
\(\overline{a_na_0}\) in slot \(0\), that is, it lets \(a_n\) act on the
left of \(W=(\fA/J)^{\hot s}\), through the first factor. This action
kills \(J\). Moving \((\bar a_0,\dots,\bar a_{s-1})\) to the end
identifies \(H^s/H^{s-1}\) with the complex \(P(W)\) of
Lemma~\ref{ha:lem:barW}, of degree \(n-s\), and the differential with
\((-1)^sb'_W\). (For \(n=s\) the cyclic face gives \(a_sa_0\in J\) in slot
\(0\), which lies in \(H^0\subset H^{s-1}\).) By Lemma~\ref{ha:lem:barW} each
quotient is acyclic. The filtration is finite in each degree, so the
inclusion \(H^0\subset H^n\) is a quasi-isomorphism near degree \(n\).
\end{proof}

\begin{theorem}[Wodzicki excision, (B4)]\label{ha:thm:excision}
Let \(J\subset\fA\) be a closed ideal of a filtered nuclear algebra, and
suppose that \(J\) is topologically H-unital. Then \(C(J)\to\cK\) is a
quasi-isomorphism, and there is a long exact sequence
\[
 \cdots\to\HH_k(J)\to\HH_k(\fA)\to\HH_k(\fA/J)\xrightarrow{\ \partial\ }
 \HH_{k-1}(J)\to\cdots .
\]
To compute \(\partial[c]\), lift the cycle \(c\) to a chain \(\tilde c\)
of \(\fA\) (Lemma~\ref{ha:lem:exact}(1)); then \(b\tilde c\in\cK\) is a
cycle, and \(\partial[c]\) is its class in \(H(\cK)\cong\HH(J)\).
\end{theorem}

\begin{proof}
Propositions~\ref{ha:prop:stepA} and~\ref{ha:prop:stepB} give the
quasi-isomorphism. Lemma~\ref{ha:lem:exact}(1) gives the short exact
sequence \eqref{eq:chainses}. Its long exact sequence, with \(H(\cK)\)
replaced by \(\HH(J)\), is the stated one, and the description of
\(\partial\) is the snake lemma.
\end{proof}

This is Background Fact~\ref{bg:excision}. Lemma~\ref{lem:conn} on
cocycles follows from it as shown there. The theorem uses neither a unit
in \(\fA\) nor a section of \(\fA\to\fA/J\). For algebras with algebraic
tensor products the same proof applies, with (F1) replaced by the
exactness of \(\otimes\) over \(\C\).

\subsection{Connes' theorem (B5)}\label{ha:subsec:connes}

We adapt \cite[Ch.~10, ``Smooth functions'']{MelroseMicro}, and add what
is needed for boundaries and corners. Let \(M\) be a compact manifold with
corners of dimension \(n\). Its charts are open subsets of
\([0,\infty)^l\times\R^{n-l}\). By (F2), \(C_k=\CI(M^{k+1})\), and the face
maps are restrictions:
\[
 (d_iu)(z_0,\dots,z_{k-1})=u(z_0,\dots,z_i,z_i,\dots,z_{k-1}),\qquad
 (d_ku)(z_0,\dots,z_{k-1})=u(z_0,\dots,z_{k-1},z_0).
\]
The HKR map is, in local coordinates,
\[
 \chi(u)(x)=\frac1{k!}\sum_{j_1,\dots,j_k}
 \frac{\partial^ku}{\partial z_1^{j_1}\cdots\partial z_k^{j_k}}(x,\dots,x)\,
 dx^{j_1}\wedge\dots\wedge dx^{j_k},
\]
with one-sided derivatives at the boundary. On elementary tensors
\(\chi(a_0\otimes\dots\otimes a_k)=\frac1{k!}a_0\,da_1\wedge\dots\wedge
da_k\). Put \(\varphi(a_0;a_1,\dots,a_k)=\sum_{\sigma\in\mathfrak S_k}
\sgn(\sigma)\,a_0\otimes a_{\sigma(1)}\otimes\dots\otimes a_{\sigma(k)}\).

\begin{lemma}\label{ha:lem:chi}
\(\chi\circ b=0\), \(b\varphi(a_0;a_1,\dots,a_k)=0\) and
\(\chi\varphi(a_0;a_1,\dots,a_k)=a_0\,da_1\wedge\dots\wedge da_k\). In
particular \(\chi_*:\HH_k(\CI(M))\to\Om^k(M)\) is onto.
\end{lemma}

\begin{proof}
All three identities are algebraic for elementary tensors. Put
\(T_i=a_0a_i\,da_1\wedge\dots\widehat{da_i}\dots\wedge da_k\). By the Leibniz
rule, \((k-1)!\,\chi(b(a_0\otimes\dots\otimes a_k))=T_1+\sum_{i=1}^{k-1}
(-1)^i(T_i+T_{i+1})+(-1)^kT_k=0\). In \(b\varphi\) the inner faces cancel in
pairs (\(\sigma\) against \(\sigma\circ(i\ i{+}1)\)), and \(d_0\) for
\(\sigma\) cancels \((-1)^kd_k\) for the cyclically shifted permutation,
whose sign is \((-1)^{k-1}\sgn\sigma\). The third identity is clear. Since
\(\chi\circ b\) is continuous and the algebraic tensor product is dense,
\(\chi\circ b=0\) on all chains. Every \(k\)-form is a finite sum
\(\sum a_0\,df_1\wedge\dots\wedge df_k\): use a partition of unity
\(\sum\rho_\beta=1\) subordinate to charts and global functions
\(f^\beta_i\) equal to the coordinates of the chart on
\(\supp\rho_\beta\).
\end{proof}

\begin{theorem}[HKR for polynomials \cite{MelroseMicro,Loday}]
\label{ha:thm:poly}
For \(\fA=\C[x_1,\dots,x_n]\), with algebraic tensor products, \(\chi\)
induces \(\HH_k(\fA)\cong\fA\otimes\Lambda^k\C^n\), with inverse
\(a_0\,dx^{i_1}\wedge\dots\wedge dx^{i_k}\mapsto\varphi(a_0;x_{i_1},
\dots,x_{i_k})\). The same holds for \(\C[r,r^{-1}]\), with
\(\HH_0=\C[r^{\pm1}]\), \(\HH_1=\C[r^{\pm1}]\,dr\), \(\HH_k=0\) for \(k\ge2\).
\end{theorem}

\begin{proof}
\(\HH_k(\fA)=\operatorname{Tor}^{\fA^e}_k(\fA,\fA)\), computed by the bar
resolution \(P_k=\fA^{\otimes(k+2)}\), which is contracted by \(1\otimes\).
The Koszul complex \(K_k=\C[x,y]\otimes\Lambda^k\C^n\), with
\(\partial e_i=x_i-y_i\), is a second free resolution. In the variables
\(u=x-y\), \(v=y\), \(\partial\) is contraction with the Euler field
\(\sum u_i\partial_{u_i}\), so the exterior derivative in \(u\) divided by
the degree in \(u\) contracts \(K\) in positive \(u\)-degree; the
augmentation is \(f\mapsto f(x,x)\). The map
\(e_{i_1}\wedge\dots\wedge e_{i_k}\mapsto\sum_\sigma\sgn\sigma\,1\otimes
x_{i_{\sigma(1)}}\otimes\dots\otimes x_{i_{\sigma(k)}}\otimes1\) is a
chain map \(K\to P\) over the identity: the inner faces cancel in pairs,
and the outer faces give \(\sum_l(-1)^{l-1}(x_{i_l}-y_{i_l})\) times the
smaller antisymmetrization. By the comparison theorem it induces an
isomorphism after \(\fA\otimes_{\fA^e}\), where the Koszul side has zero
differential. This gives the inverse map \(\varphi\), and
\(\chi\varphi=1\) by Lemma~\ref{ha:lem:chi}. For \(\C[r^{\pm1}]\) the
resolution is \(0\to\fA^e\xrightarrow{r-s}\fA^e\to\fA\to0\). It is exact
because \(r-s\) is not a zero divisor and \(f(r,s)-f(r,r)\) is divisible
by \(r-s\) (clear the denominator \((rs)^N\) first).
\end{proof}

\subsubsection*{Normalized chains}
The next three results are algebraic and hold verbatim for \(\CI(M)\), for
the polynomial algebras, and for the formal and Laurent algebras of
Theorem~\ref{ha:thm:variants}. Let \(D_i=\{z_i=z_{i+1}\}\subset M^{k+1}\)
for \(0\le i<k\) and \(D_k=\{z_k=z_0\}\), and
\begin{gather*}
 N_k=\{u\in C_k:\ u|_{D_i}=0,\ 0\le i<k\},\qquad M_k=\{u\in N_k:\
 u|_{D_k}=0\}\quad(k\ge1),\\
 B_k=d_{k+1}(N_{k+1}),\qquad N_0=M_0=C_0 .
\end{gather*}
If \(u\in N_k\) then \(bu=(-1)^kd_ku\), and \(d_ku\in N_{k-1}\). So \(N\) is
a subcomplex. Moreover \(B_k\subset M_k\), and \(H_k(N)=M_k/B_k\).

\begin{proposition}[{\cite[Ch.~10]{MelroseMicro}}]\label{ha:prop:normalized}
\(N\subset C\) is a quasi-isomorphism. Hence \(\HH_k=M_k/B_k\), and an
element of \(M_k\) that is a boundary in \(C\) lies in \(B_k\).
\end{proposition}

\begin{proof}
Let \(J^j_m=\{u\in C_m:\ u|_{D_i}=0\text{ for }i<\min(j,m)\}\); each
\(J^j\) is a subcomplex, \(J^0=C\), and \(J^j_m=N_m\) for \(m\le j\). For
\(m\ge j+1\), every \(u\in J^j_m\) splits uniquely as \(u=u'+u''\), where
\(u'(z)=u(z_0,\dots,z_j,z_j,z_{j+2},\dots,z_m)\) does not depend on
\(z_{j+1}\) and \(u''=u-u'\in J^{j+1}_m\). Let \(V^j_m\) be the
\(u\in J^j_m\) that do not depend on \(z_{j+1}\) (and \(V^j_m=0\) for
\(m\le j\)). Then \(J^j=J^{j+1}\oplus V^j\) as graded spaces. On \(V^j\)
the faces \(d_i\), \(i<j\), vanish, and \(d_j=d_{j+1}\), so
\(bu=\sum_{i\ge j+2}(-1)^id_iu\in V^j\). The map
\((Hu)(z_0,\dots,z_{m+1})=u(z_0,\dots,z_{j+1},z_{j+3},\dots,z_{m+1})\),
which inserts a dummy variable in position \(j+2\), satisfies
\(d_{j+2}H=1\) and \(d_iH=Hd_{i-1}\) for \(i\ge j+3\). Hence
\(bH+Hb=(-1)^j\) on \(V^j\), so \(V^j\) is a contractible subcomplex, and
\(J^{j+1}\subset J^j\) is a quasi-isomorphism. Therefore \(H_k(C)\cong
H_k(J^{k+1})=H_k(N)\).
\end{proof}

\begin{lemma}\label{ha:lem:Bmodule}
\textup{(i)} \(B_k\) is a \(\CI(M^{k+1})\)-submodule of \(M_k\).
\textup{(ii)} If \(f\in\CI(M)\) and \(g\in M_k\), then
\((f(z_j)-f(z_{j+1}))g\in B_k\) for \(0\le j<k\), and
\((f(z_k)-f(z_0))g\in B_k\).
\end{lemma}

\begin{proof}
(i) If \(v\in N_{k+1}\) and \(\psi\in\CI(M^{k+1})\), then
\(\tilde\psi v\in N_{k+1}\) with \(\tilde\psi(z_0,\dots,z_{k+1})=
\psi(z_0,\dots,z_k)\), and \(d_{k+1}(\tilde\psi v)=\psi\,d_{k+1}v\). (ii)
For \(j<k\) let \(w=(-1)^{j+1}(f(z_j)-f(z_{j+1}))\,g(z_0,\dots,z_j,
z_{j+2},\dots,z_{k+1})\in C_{k+1}\). All faces of \(w\) vanish except
\(d_{j+1}\), because \(g\in M_k\) and because of the first factor. So
\(bw=(f(z_j)-f(z_{j+1}))g\). This lies in \(M_k\) and is a boundary, hence
in \(B_k\) by Proposition~\ref{ha:prop:normalized}. The last case is the
sum of the others.
\end{proof}

\begin{example}[The model \(k=1\)]\label{ha:ex:k1}
Let \(M=[0,1]\) and \(u\in M_1\), so \(u(z_0,z_0)=0\) and
\(u=(z_1-z_0)w\). Then \(\chi(u)=w(x,x)\,dx\). If \(\chi(u)=0\), then
\(u=(z_1-z_0)^2g\) with \(g\) smooth, and \(u=d_2v\) for
\(v(z_0,z_1,z_2)=-(z_1-z_0)(z_2-z_1)g(z_0,z_1)\in N_2\). So \(u\in B_1\), and
\(\HH_1(\CI([0,1]))=\Om^1([0,1])\). The proof below does the same in
general: Taylor expansion at the diagonal, a polynomial lift, and
division of a flat remainder.
\end{example}

\subsubsection*{The local structure of normalized cycles}
Near a point \((x_0,\dots,x_0)\) of the diagonal \(\Delta\) choose a chart
around \(x_0\) whose image \(U\) is a \emph{convex} subset of
\([0,\infty)^l\times\R^{n-l}\), and use the coordinates
\[
 x=z_0,\qquad y_j=z_j-z_{j-1}\quad(1\le j\le k)
\]
on \(U^{k+1}\). Then \(D_{j-1}=\{y_j=0\}\), \(D_k=\{y_1+\dots+y_k=0\}\) and
\(\Delta=\{y=0\}\). The set \(U^{k+1}\) is convex with non-empty interior,
so the Taylor series in \(y\) at points of \(\Delta\) of a smooth function
on it is determined by the function. It vanishes on a linear subspace
\(S\) if the function vanishes on \(S\cap U^{k+1}\), because \(S\cap
U^{k+1}\) is convex with non-empty interior in \(S\) and contains
\(\Delta\cap U^{k+1}\). This replaces the use of Hadamard's lemma in
\cite{MelroseMicro}, whose integration paths can leave a half-space.

\begin{lemma}\label{ha:lem:antisym}
A \(k\)-linear form \(q\) on \(\R^n\) satisfies
\(q(y_1,\dots,y_{k-1},-(y_1+\dots+y_{k-1}))\equiv0\) if and only if it is
alternating.
\end{lemma}

\begin{proof}
If \(q\) is alternating the expression is \(-\sum_{j<k}q(y_1,\dots,
y_{k-1},y_j)=0\). Conversely, the terms \(q(y_1,\dots,y_{k-1},y_j)\) have
distinct multidegrees, so each vanishes. Thus \(q\) is alternating in the
arguments \(j,k\), and these transpositions generate \(\mathfrak S_k\).
\end{proof}

\begin{lemma}\label{ha:lem:leading}
Let \(u\in N_k\). The Taylor series of \(u\) at \(\Delta\) in \(y\) has
coefficients smooth in \(x\in U\). Each of its monomials contains a
component of every \(y_j\). Its part of degree \(k\) is
\(u_{(k)}=\sum_Iw_I(x)\,y_{1,i_1}\cdots y_{k,i_k}\), and
\(\chi(u)=\frac1{k!}\sum_Iw_I\,dx^{i_1}\wedge\dots\wedge dx^{i_k}\). If
\(u\in M_k\), then \(u_{(k)}(x;\cdot)\) is alternating, and \(\chi(u)(x)=0\)
if and only if \(u_{(k)}(x;\cdot)=0\).
\end{lemma}

\begin{proof}
The Taylor series vanishes on each \(\{y_j=0\}\), which gives the first
claim. Since \(z_l=x+y_1+\dots+y_l\), we have \(\partial_{z_l}=
\partial_{y_l}-\partial_{y_{l+1}}\) (and \(\partial_{z_k}=\partial_{y_k}\)).
In \(\partial_{z_1^{i_1}}\cdots\partial_{z_k^{i_k}}u\) at \(y=0\) only the
term in which each \(y_l\) is differentiated once survives, and it equals
\(w_I(x)\). If \(u\in M_k\), then \(u_{(k)}\) vanishes on
\(\{y_1+\dots+y_k=0\}\); apply Lemma~\ref{ha:lem:antisym}.
\end{proof}

\begin{lemma}[{\cite[Ch.~10]{MelroseMicro}}]\label{ha:lem:polylift}
Let \(k\ge1\), \(m\ge k+1\). Let \(\mathcal M_m\) be the homogeneous
polynomials of degree \(m\) in \((y_1,\dots,y_k)\) vanishing on each
\(\{y_j=0\}\) and on \(\{y_1+\dots+y_k=0\}\). Let \(\mathcal V_m\) be the
homogeneous polynomials of degree \(m\) in \((y_1,\dots,y_{k+1})\)
vanishing on each \(\{y_j=0\}\). Then \(R:\mathcal V_m\to\mathcal M_m\),
\((RV)(y_1,\dots,y_k)=V(y_1,\dots,y_k,-(y_1+\dots+y_k))\), is onto.
\end{lemma}

\begin{proof}
Given \(q\in\mathcal M_m\), the polynomial chain \(u(z)=q(z_1-z_0,\dots,
z_k-z_{k-1})\) lies in \(M_k\) and \(\chi(u)=0\). By
Theorem~\ref{ha:thm:poly} and Proposition~\ref{ha:prop:normalized} for
polynomials, \(u=d_{k+1}v\) with \(v\in N_{k+1}\). Write \(v\) in the
coordinates \(x,y_1,\dots,y_{k+1}\). Its component of degree \(0\) in
\(x\) and \(m\) in \(y\) lies in \(\mathcal V_m\), and \(R\) maps it to
\(q\).
\end{proof}

\begin{lemma}\label{ha:lem:local}
Let \(k\ge1\) and \(u\in M_k\) with \(\chi(u)=0\). Every point \(p\in
M^{k+1}\) has a neighbourhood on which \(u\) agrees with an element of
\(B_k\).
\end{lemma}

\begin{proof}
If \(p\notin\Delta\), then \(p_j\ne p_{j+1}\) for some \(j<k\). Choose
\(f\in\CI(M)\) with \(f(p_j)\ne f(p_{j+1})\), and \(\psi\in\CI(M^{k+1})\)
equal to \(1\) near \(p\) and supported where \(f(z_j)\ne f(z_{j+1})\).
Then \(g=\psi u/(f(z_j)-f(z_{j+1}))\in M_k\), and near \(p\),
\(u=(f(z_j)-f(z_{j+1}))g\in B_k\) by Lemma~\ref{ha:lem:Bmodule}.

Let \(p=(x_0,\dots,x_0)\). By Lemma~\ref{ha:lem:leading} and \(\chi(u)=0\),
the Taylor series of \(u\) at \(\Delta\) is \(\sum_{m\ge k+1}u_{(m)}\) with
\(u_{(m)}\in\CI(U)\otimes\mathcal M_m\). Using a linear right inverse of
\(R\), choose \(V_{(m)}\in\CI(U)\otimes\mathcal V_m\) with
\(RV_{(m)}=u_{(m)}\). By Borel's lemma, \(V=\sum_m\chi_0(\lambda_m|y|)
V_{(m)}\), with \(\chi_0\in\CI_c(\R)\) equal to \(1\) near \(0\) and
\(\lambda_m\to\infty\) fast enough, is smooth for \(x\in U\) up to the
boundary and small \(y\in(\R^n)^{k+1}\). It has Taylor series
\(\sum V_{(m)}\) and vanishes on each \(\{y_j=0\}\). Let \(\Theta\) be a
cut-off near \((x_0,\dots,x_0)\in M^{k+2}\), and let \(v=\Theta V\),
written in \(z\). Then \(v\in N_{k+1}\). The chain \(s=d_{k+1}v\in B_k\)
has the same Taylor series as \(u\) at \(\Delta\) near \(p\), since
\(z_{k+1}=z_0\) means \(y_{k+1}=-(y_1+\dots+y_k)\).

So \(g=u-s\in M_k\) vanishes to infinite order at \(\Delta\) near \(p\).
For \(z\in U^{k+1}\) the segment from \((z_0,\dots,z_0)\) to \(z\) lies in
\(U^{k+1}\), by convexity. Taylor's formula along it gives
\(|\partial^\alpha g(z)|\le C_{\alpha,N}|y|^N\) for all \(N\). Hence
\(g_{j,i}=y_{j,i}g/|y|^2\), extended by \(0\) on \(\Delta\), is smooth up
to the boundary, and it vanishes where \(g\) does. Let \(f_i\in\CI(M)\)
equal the \(i\)-th coordinate near \(x_0\), and let \(\psi\) be a cut-off
equal to \(1\) near \(p\). Then near \(p\),
\(g=\sum_{j,i}y_{j,i}g_{j,i}=\sum_{j,i}(f_i(z_j)-f_i(z_{j-1}))\psi
g_{j,i}\), which lies in \(B_k\) by Lemma~\ref{ha:lem:Bmodule}. So
\(u=s+g\) agrees with an element of \(B_k\) near \(p\).
\end{proof}

\begin{theorem}[Connes, (B5)]\label{ha:thm:connes}
For a compact manifold with corners \(M\), \(\chi\) induces isomorphisms
\(\HH_k(\CI(M))\cong\Om^k(M)\).
\end{theorem}

\begin{proof}
For \(k=0\), \(b=0\) on \(C_1\). For \(k\ge1\), \(\chi_*\) is onto by
Lemma~\ref{ha:lem:chi}. By Proposition~\ref{ha:prop:normalized} a class is
represented by \(u\in M_k\), and it vanishes if and only if \(u\in B_k\).
If \(\chi(u)=0\), Lemma~\ref{ha:lem:local} and compactness give a finite
partition of unity \(\sum\psi_\beta=1\) on \(M^{k+1}\) and
\(s_\beta\in B_k\) with \(u=s_\beta\) on \(\supp\psi_\beta\). Then
\(u=\sum_\beta\psi_\beta s_\beta\in B_k\) by Lemma~\ref{ha:lem:Bmodule}(i).
\end{proof}

\subsubsection*{The coefficient algebras of the notes}
Let \(N\) be a compact manifold with corners and \(x\) a formal variable.
Let \(M\) be a compact manifold with boundary, with boundary defining
function \(x\) and a collar \([0,\epsilon)\times\partial M\). We consider
\[
 \CI(N)[[x]],\qquad \CI(N)\Lser{x},\qquad x^{-\Z}\CI(M),\qquad\dCI(M).
\]
The localized algebras are filtered nuclear with pieces \(x^{-L}\CI(N)
[[x]]\) and \(x^{-L}\CI(M)\). By (F2) and (F3) their chains are
\(\CI(N^{k+1})[[x_0,\dots,x_k]]\), the unions over \(L\) of
\((x_0\cdots x_k)^{-L}\) times the previous chains, and \(\dCI(M^{k+1})\).
The forms are \(\Om(N)[[x]]\oplus\Om(N)[[x]]\,dx\) (with \(d\) the formal
differential), the corresponding Laurent forms, \(x^{-\Z}\Om(M)\), and the
flat forms \(\dOm(M)\). The HKR map is given by the same formula.

\begin{theorem}[HKR for the coefficient algebras, (B5)]
\label{ha:thm:variants}
For each of the four algebras \(\chi\) induces an isomorphism of
\(\HH_k\) onto the \(k\)-forms just listed. The isomorphisms commute with
Taylor expansion at \(\partial M\) in the collar. The same holds for
\(\cS(\R^d)\), whose forms are the Schwartz forms.
\end{theorem}

\begin{proof}
\emph{Formal series.} Run the proof of Theorem~\ref{ha:thm:connes} with
\(x\) as an extra coordinate in which everything is a formal power series.
Normalized chains and Lemma~\ref{ha:lem:Bmodule} are algebraic.
Differences \(x_j-x_{j-1}\) are formal variables, and a series in
\(x_0,\dots,x_k\) is a series in \(x_0\) and the differences. Points of
\(N^{k+1}\) off the \(N\)-diagonal are treated as before. Near the
\(N\)-diagonal, expand in all differences, lift with
Lemma~\ref{ha:lem:polylift} (polynomials in the \(N\)-coordinates and
\(x\)), and apply Borel's lemma only in the \(N\)-differences. For each
fixed monomial in \(x_0\) and the \(x\)-differences this is a convergent
Borel sum; choose \(\lambda_m\) to control all monomials of degree \(\le
m\). Every coefficient of the remainder \(g\) is then flat at the
\(N\)-diagonal, so it can be divided by \(|y_N|^2\) coefficientwise. If
\(\dim N=0\) the remainder is \(0\).

\emph{Laurent coefficients.} Let \(R\) be \(\CI(M)\) or \(\CI(N)[[x]]\)
and \(R'=R[x^{-1}]\). Proposition~\ref{ha:prop:normalized} holds for
\(R'\) as stated. Let \(u\in M_k(R')\) with \(\chi(u)=0\), and write
\(u=fv\) with \(f=(x_0\cdots x_k)^{-L}\) and \(v\in C_k(R)\). Multiplication
by \(f\) is injective, so \(v\in M_k(R)\). By Lemma~\ref{ha:lem:leading},
\(v=O(|y|^k)\) at \(\Delta\), so the degree-\(k\) Taylor part of \(fv\) is
\(f|_\Delta v_{(k)}\) and \(\chi(u)=x^{-(k+1)L}\chi(v)\) (as an identity of
forms on the interior, hence of Laurent forms). So \(\chi(v)=0\), hence
\(v=d_{k+1}w\) with \(w\in N_{k+1}(R)\), and \(u=d_{k+1}(\tilde fw)\in
B_k(R')\), where \(\tilde f(z_0,\dots,z_{k+1})=f(z_0,\dots,z_k)\).
Surjectivity is Lemma~\ref{ha:lem:chi} with \(a_0\in x^{-L}R\).

\emph{Flat functions.} Taylor expansion in the collar gives the exact
sequence \(0\to\dCI(M)\to\CI(M)\to\CI(\partial M)[[x]]\to0\) (Borel's
lemma). Its kernel is topologically H-unital by
Lemma~\ref{lem:flatHunital}, so Theorem~\ref{ha:thm:excision} gives a
long exact sequence. The HKR maps commute with Taylor expansion, since
both sides are computed from derivatives at the diagonal. The sequence of
forms \(0\to\dOm\to\Om(M)\to\Om(\partial M)[[x]]\oplus\Om(\partial M)
[[x]]\,dx\to0\) is exact, by Borel's lemma for forms. The HKR maps for the
middle and right algebras are isomorphisms, so \(\HH_k(\CI(M))\to
\HH_k(\CI(\partial M)[[x]])\) is onto for every \(k\). Hence the connecting
maps vanish, \(\HH_k(\dCI(M))\) is the kernel, and \(\chi\) identifies it
with \(\dOm^k(M)\). Finally \(\cS(\R^d)=\dCI(\overline{\R^d})\) for the
radial compactification, a closed ball, and (F2) identifies the chains.
\end{proof}

Theorem~\ref{ha:thm:variants} is the scalar case of \cite[Lemma
``Continuous HKR for the coefficient kernels'']{HHb}, which proves the flat
case directly with a Koszul complex. The version with an adjoined
polynomial or Laurent polynomial variable is
Corollary~\ref{ha:cor:homog} below.

\subsection{Künneth and Morita (B6)}\label{ha:subsec:B6}

\subsubsection*{The Künneth formula}
For unital algebras \(\fA,\mathfrak B\) the \emph{shuffle map} is
\begin{multline*}
 \mathrm{sh}\bigl((a_0\otimes\dots\otimes a_p)\otimes(b_0\otimes\dots
 \otimes b_q)\bigr)\\
 =\sum_{\sigma\in\Sh(p,q)}\sgn(\sigma)\,(a_0\otimes b_0)
 \otimes\sigma\cdot\bigl(a_1\otimes1,\dots,a_p\otimes1,1\otimes b_1,\dots,
 1\otimes b_q\bigr),
\end{multline*}
where \(\sigma\) runs over the \((p,q)\)-shuffles: it puts the entries
\(a_i\otimes1\) in positions \(\sigma(1)<\dots<\sigma(p)\) and the
\(1\otimes b_j\) in order in the remaining positions.

\begin{background}[Eilenberg--Zilber (quoted: \cite{Loday}, \S4.2)]
\label{ha:bg:EZ}
\(\mathrm{sh}:\operatorname{Tot}(C(\fA)\otimes C(\mathfrak B))\to
C(\fA\otimes\mathfrak B)\) is a chain map and a chain homotopy
equivalence. Its homotopy inverse (Alexander--Whitney) and both homotopies
are given by universal finite formulas in the face maps and in the
degeneracies \(s_i\) (insertion of \(1\) after slot \(i\)).
\end{background}

Because the formulas are universal and finite, the same holds for unital
filtered nuclear algebras with completed tensor products, using
\(C_k(\fA\hot\mathfrak B)=\fA^{\hot(k+1)}\hot\mathfrak B^{\hot(k+1)}\).

\begin{proposition}[Künneth, (B6)]\label{ha:prop:kunneth}
Let \(\fA\) be a unital filtered nuclear algebra and \(\mathfrak B\) a
unital algebra of countable dimension. Put \(C_k(\fA\otimes\mathfrak B)=
\fA^{\hot(k+1)}\otimes\mathfrak B^{\otimes(k+1)}\). Then
\[
 \HH_n(\fA\otimes\mathfrak B)\cong\bigoplus_{p+q=n}\HH_p(\fA)\otimes
 \HH_q(\mathfrak B),
\]
induced by \(\mathrm{sh}\). If both algebras are commutative, then
\(\chi(\mathrm{sh}(\alpha\otimes\beta))=\chi(\alpha)\wedge\chi(\beta)\).
\end{proposition}

\begin{proof}
By Background Fact~\ref{ha:bg:EZ}, \(\mathrm{sh}\) is a quasi-isomorphism.
The complex \(C(\fA)\otimes C(\mathfrak B)\) is an algebraic tensor product
of complexes over \(\C\), so its homology is \(H(C(\fA))\otimes
H(C(\mathfrak B))\). For the last statement, each of the
\(\binom{p+q}p\) shuffles contributes
\(a_0b_0\,da_1\wedge\dots\wedge da_p\wedge db_1\wedge\dots\wedge db_q\),
since the sign of the reordering cancels \(\sgn\sigma\). Multiplying by
\(1/(p+q)!\) gives \(\frac1{p!q!}(a_0\,da_1\wedge\cdots)\wedge(b_0\,db_1
\wedge\cdots)\).
\end{proof}

For two nuclear Fr\'echet algebras the topological Künneth formula needs
a further condition on the boundaries (quoted: \cite{Karoubi}); the notes
do not need it. What they need is the following.

\begin{corollary}[HKR with a radial variable]\label{ha:cor:homog}
Let \(R\) be one of the algebras of Theorems~\ref{ha:thm:connes}
and~\ref{ha:thm:variants}, and let \(\mathfrak B=\C[r,r^{-1}]\) (\(r\) of
weight \(1\)) or \(\C[t]\) (\(t\) of weight \(-1\)). Grade
\(C(R\otimes\mathfrak B)\) by total weight. Then \(\chi\) induces, in each
weight \(w\),
\[
 \HH_k(R\otimes\mathfrak B)_{(w)}\cong\bigl(\Om^k(R)\otimes\mathfrak B
 \oplus\Om^{k-1}(R)\otimes\mathfrak B\,d\varrho\bigr)_{(w)},\qquad
 \varrho\in\{r,t\},
\]
where \(d\varrho\) has the weight of \(\varrho\). In particular, if \(V\to
Y\) is a vector bundle with a fibre norm and \(R=\CI(SV)\) (or its Laurent,
formal or flat version), the chains of total homogeneity \(p\) have
homology the homogeneous \(k\)-forms of degree \(p\) on \(V\setminus0\).
\end{corollary}

\begin{proof}
For unital \(R\) combine Proposition~\ref{ha:prop:kunneth} with
Theorems~\ref{ha:thm:poly}, \ref{ha:thm:connes}
and~\ref{ha:thm:variants}. Every map preserves the weight. For
\(\dCI(M)\otimes\mathfrak B\) use the Taylor sequence tensored with
\(\mathfrak B\) and argue as in the flat case of
Theorem~\ref{ha:thm:variants}. The kernel has the extra degeneracy
\(E\otimes(1\otimes\cdot)\), the tensor product of extra degeneracies. A
homogeneous form of degree \(p\) on \(V\setminus0\cong SV\times(0,\infty)
_r\) is \(r^p(\alpha+\beta\wedge dr/r)\) with \(\alpha,\beta\) forms on
\(SV\), and a homogeneous function of degree \(p\) is \(r^p\) times a
function on \(SV\).
\end{proof}

\subsubsection*{Morita invariance}

\begin{lemma}[Morita for smoothing kernels; adapted from {\cite[Lemma
``Morita invariance for smoothing kernels'']{HHb}}]\label{ha:lem:morita}
Let \(\mathcal K=V\hot W\) be as in Lemma~\ref{lem:Hunital}, with
\(\langle\psi,\tilde\psi\rangle=1\), \(e=\tilde\psi\otimes\psi\) and
\(\Tr(v\otimes w)=\langle w,v\rangle\). Let \(\fA\) be unital or
topologically H-unital. Then
\begin{gather*}
 \iota(a_0\otimes\dots\otimes a_k)=(e\otimes a_0)\otimes\dots\otimes
 (e\otimes a_k),\\
 \Tr_*\bigl((\kappa_0\otimes a_0)\otimes\dots\otimes
 (\kappa_k\otimes a_k)\bigr)=\Tr(\kappa_0\cdots\kappa_k)\,a_0\otimes\dots
 \otimes a_k
\end{gather*}
are chain maps between \(C(\fA)\) and \(C(\mathcal K\hot\fA)\), with
\(\Tr_*\iota=1\), and \(\Tr_*\) is a quasi-isomorphism. For
\(V=\C^N\), \(W=(\C^N)^*\) this gives \(\HH(M_N(\fA))\cong\HH(\fA)\).
\end{lemma}

\begin{proof}
Write a chain as \((v_0\otimes w_0\otimes a_0)\otimes\dots\otimes(v_k
\otimes w_k\otimes a_k)\) and rearrange it as the \emph{letters}
\(\ell_i=w_i\otimes v_{i+1}\in U=W\hot V\) (indices mod \(k+1\)) together
with \(a_0\otimes\dots\otimes a_k\). Let \(\varepsilon(w\otimes v)=
\langle w,v\rangle\). Since \((v_i\otimes w_i)(v_{i+1}\otimes w_{i+1})=
\varepsilon(\ell_i)\,v_i\otimes w_{i+1}\), the face \(d_i\) applies
\(\varepsilon\) to \(\ell_i\), deletes it, and multiplies \(a_i\) with
\(a_{i+1}\) (for \(i=k\), \(a_ka_0\) in slot \(0\)); the other letters are
unchanged. Split \(U=\C u_0\oplus U'\) with \(u_0=\psi\otimes\tilde\psi\)
and \(U'=\ker\varepsilon\). The number \(m\) of letters in \(U'\) is
preserved by \(b\), so \(C=\bigoplus_mC^{(m)}\).

If \(m=0\), every letter is \(u_0\), every slot is \(\tilde\psi\otimes
\psi=e\), and \(C^{(0)}=\iota C(\fA)\) with the differential of \(\fA\).
The projection onto \(C^{(0)}\) replaces each letter \(\ell\) by
\(\varepsilon(\ell)u_0\), that is, it is \(\iota\Tr_*\), since
\(\prod_i\varepsilon(\ell_i)=\Tr(\kappa_0\cdots\kappa_k)\).

If \(m\ge1\), filter by the position \(w\) of the last \(U'\)-letter. A
face \(d_i\) with \(i<w\) lowers \(w\), \(d_w\) vanishes, and \(d_i\) with
\(i>w\) keeps \(w\). On the associated graded complex only the faces
\(d_{w+1},\dots,d_k\) survive. They multiply neighbours in the block
\((a_{w+1},\dots,a_k,a_0)\), of length \(k-w+1\ge1\), with signs
\(\pm(-1)^{i-w-1}\). So the graded complex is the bar complex of \(\fA\)
tensored with spectators, and it is contracted by \(1\otimes\) or by
\(E\), inserted in front of the block. The filtration is finite in each
degree, so \(C^{(m)}\) is acyclic for \(m\ge1\).
\end{proof}

For full corners we use a homotopy that we have not found written out in
this form; it is the chain-level version of \cite[\S1.2]{Loday}. Let
\(T\) be an algebra containing \(e,P,Q\) with \(e^2=e\), \(PQ=e\),
\(eP=P\) and \(Qe=Q\). Then \(a\mapsto QaP\) is an algebra map
\(eTe\to T\).

\begin{lemma}\label{ha:lem:PQ}
Let \(\phi_0,\phi_1:C(eTe)\to C(T)\) be induced by the inclusion and by
\(a\mapsto QaP\). For \(0\le j\le k\) put
\[
 h_j(a_0\otimes\dots\otimes a_k)=a_0P\otimes Qa_1P\otimes\dots\otimes
 Qa_jP\otimes Q\otimes a_{j+1}\otimes\dots\otimes a_k,
\]
and \(H=\sum_j(-1)^jh_j\). Then \(bH+Hb=\phi_0-\phi_1\).
\end{lemma}

\begin{proof}
Use \(a_iPQ=a_i\) and \(QaPQ=Qa\) for \(a,a_i\in eTe\). In \(h_j\) the
face \(d_i\) multiplies two transformed entries if \(i<j\), giving
\(h_{j-1}d_i\). If \(i>j+1\) it multiplies two original entries (or
\(a_k\) with \(a_0P\)), giving \(h_jd_{i-1}\). The remaining faces satisfy
\[
 d_0h_0=1,\qquad d_{k+1}h_k=\phi_1,\qquad d_jh_j=d_jh_{j-1}\ (j\ge1),
\]
both sides of the last being \((\dots,Qa_{j-1}P,Qa_j,a_{j+1},\dots)\). The
terms with \(i\in\{j,j+1\}\) in \(bH\) therefore telescope to
\(d_0h_0-d_{k+1}h_k=\phi_0-\phi_1\). The terms with \(i<j\) or \(i>j+1\)
add up to \(-Hb\). (We also checked the identity by computer algebra for
\(k\le4\).)
\end{proof}

\begin{proposition}[Morita for full corners, (B6)]\label{ha:prop:corner}
Let \(S\) be a unital filtered nuclear algebra and \(f\in S\) an
idempotent. Suppose there are \(u_1,\dots,u_m\in Sf\) and
\(v_1,\dots,v_m\in fS\) with \(\sum_ju_jv_j=1\). Then the inclusion
\(fSf\to S\) induces an isomorphism \(\HH(fSf)\cong\HH(S)\). Its inverse
is induced by
\(\theta(s_0\otimes\dots\otimes s_k)=\sum_{i_0,\dots,i_k}v_{i_0}s_0u_{i_1}
\otimes v_{i_1}s_1u_{i_2}\otimes\dots\otimes v_{i_k}s_ku_{i_0}\).
\end{proposition}

\begin{proof}
In \(T=M_{1+m}(S)\), with indices \(0,\dots,m\), take \(e=E_{00}\),
\(P=\sum_ju_jE_{0j}\) and \(Q=\sum_jv_jE_{j0}\). Then \(PQ=e\), \(eTe=S\),
and \(QsP=\sum_{i,j}v_isu_jE_{ij}\). By Lemma~\ref{ha:lem:PQ} and the
generalized trace \(\Tr_*:C(M_{1+m}(S))\to C(S)\) (a chain map with
\(\Tr_*\phi_0=1\)), \(1\simeq\Tr_*\phi_1=\theta\) on \(C(S)\). The image of
\(\theta\) lies in \(C(fSf)\). In \(T'=M_{1+m}(fSf)\) take \(e'=fE_{00}\),
\(P'=\sum_jfu_jE_{0j}\) and \(Q'=\sum_jv_jfE_{j0}\). All entries lie in
\(fSf\), \(P'Q'=e'\), and \(Q'aP'=\sum v_iau_jE_{ij}\) for \(a\in fSf\).
The same argument gives \(1\simeq\theta\) on \(C(fSf)\).
\end{proof}

\begin{remark}[Bundles]\label{ha:rem:bundles}
For a subbundle \(E\subset X\times\C^N\) of positive rank, with orthogonal
projection \(e\), the hypothesis holds in \(S=M_N(\CI(X))\subset
M_N(\fA)\), for every algebra \(\fA\) of the notes. Take a partition of
unity \(\sum\rho_\alpha^2=1\) and unit sections \(s_\alpha\) of \(E\) on
\(\supp\rho_\alpha\), and put \(u_{\alpha,i}=\rho_\alpha e_is_\alpha^*\) and
\(v_{\alpha,i}=\rho_\alpha s_\alpha e_i^*\). Then \(u_{\alpha,i}e=
u_{\alpha,i}\), \(ev_{\alpha,i}=v_{\alpha,i}\) and \(\sum_{\alpha,i}
u_{\alpha,i}v_{\alpha,i}=1\). With Lemma~\ref{ha:lem:morita} this gives
\(\HH(eM_N(\fA)e)\cong\HH(\fA)\) for the unital algebras of the notes, as
used in Section~\ref{sec:morita}. For the ideals \(J\) (\(\cI\), \(\Is\),
\(\Id\)) the five lemma and Theorem~\ref{ha:thm:excision} give the same,
provided \(eM_N(J)e\) is topologically H-unital. This holds because the
contractions of Lemmas~\ref{lem:Hunital}, \ref{lem:flatHunital} and
Propositions~\ref{prop:Is}, \ref{prop:Id} are defined by kernels and can
be written with bundle coefficients (sketch).
\end{remark}

\subsection{Convergence of the order spectral sequence}
\label{ha:subsec:conv}

\subsubsection*{The algebraic statement}
Let \((C_\bullet,b)\) be a complex in which each \(C_k\) is a subspace of
the space of series \(c=\sum_{p\le P}c_p\), \(c_p\in C_{p,k}\), containing
each \(C_{p,k}\), with \(b=\sum_{l\ge0}b_l\) and \(b_l:C_{p,k}\to
C_{p-l,k-1}\). Put \(F_pC=\{c:\ c_q=0\ \text{for}\ q>p\}\),
\(E^1_{p,k}=H_k(C_{p,\bullet},b_0)\), and let \(d^1\) be induced by \(b_1\).
Comparing degrees in \(b^2=0\) gives \(\sum_{i+j=l}b_ib_j=0\), so \(d^1\) is
well defined, \((d^1)^2=0\), and \(E^2=H(E^1,d^1)\). The filtration is
\emph{complete} if every series \(\sum_je_j\) with \(e_j\in F_{p_j}C_k\),
\(\sup p_j<\infty\) and \(p_j\to-\infty\) defines an element of \(C_k\), and
\(b\) commutes with such sums.

\begin{proposition}[Algebraic convergence {\cite[Ch.~10]{MelroseMicro}}]
\label{ha:prop:algconv}
Suppose that, for an integer \(n\), \(E^2_{p,k}=0\) whenever \(p\ne k-n\).
\begin{enumerate}
\item Every class in \(H_k(C,b)\) has a representative in \(F_{k-n}C_k\),
and \(\Theta[c]=[c_{k-n}]\in E^2_{k-n,k}\) is well defined.
\item \(\ker\Theta=\bigcap_pF_pH_k\), where \(F_pH_k\) is the image of
\(H_k(F_pC)\).
\end{enumerate}
If the filtration is complete, then \(\Theta:H_k(C,b)\to E^2_{k-n,k}\) is an
isomorphism, and all differentials \(d^r\), \(r\ge2\), vanish.
\end{proposition}

\begin{proof}
\emph{Step (P).} Let \(c\in F_pC_k\) with \(b_0c_p=0\) and \(b_1c_p+b_0
c_{p-1}=0\), and suppose \([c_p]=0\) in \(E^2_{p,k}\). Then \([c_p]=d^1[x]\)
for a \(b_0\)-cycle \(x\in C_{p+1,k+1}\), so \(c_p-b_1x=b_0y\) with
\(y\in C_{p,k+1}\). The components of \(b(x+y)\) of order \(p+1\) and \(p\)
are \(b_0x=0\) and \(b_1x+b_0y=c_p\). Hence \(c-b(x+y)\in F_{p-1}\).

\emph{(1).} If \(c\in F_PC_k\) is a cycle and \(P>k-n\), then
\(E^2_{P,k}=0\), so (P) lowers \(P\); after finitely many steps \(c\in
F_{k-n}\). If \(c\in F_{k-n}\) is a cycle, then \(b_0c_{k-n}=0\) and
\(b_1c_{k-n}=-b_0c_{k-n-1}\), so \(c_{k-n}\) defines a class in \(E^2\). Let
\(c=be\in F_{k-n}\) with \(e\in F_q\). If \(q>k+1-n\), then (P) applies
to \(e\), since \(b_0e_q=(be)_q=0\), \(b_1e_q+b_0e_{q-1}=(be)_{q-1}=0\) and
\(E^2_{q,k+1}=0\). It replaces \(e\) by an element of \(F_{q-1}\) with the
same boundary. Once \(q\le k+1-n\), \(c_{k-n}=b_1e_{k+1-n}+b_0e_{k-n}\), so
\([c_{k-n}]=d^1[e_{k+1-n}]=0\) in \(E^2\).

\emph{(2).} If \([c]\in F_{k-n-1}H_k\), then \(\Theta[c]=0\) by (1).
Conversely, if \(c\in F_{k-n}\) and \([c_{k-n}]=0\), then (P) gives a
homologous cycle in \(F_{k-n-1}\). Since \(E^2_{p,k}=0\) for \(p<k-n\),
repeating (P) finitely often reaches any \(F_p\).

\emph{Injectivity.} In the situation of (2), the repeated steps give
\(e_p\in F_{p+1}C_{k+1}\), \(p=k-n,k-n-1,\dots\), with \(c-b\sum_{p\ge
P}e_p\in F_{P-1}\). By completeness \(c=b\sum_pe_p\).

\emph{Surjectivity.} Let \(x\in C_{k-n,k}\) be a \(b_0\)-cycle with
\(d^1[x]=0\). Then \(b_1x=-b_0y\), and \(c^{(1)}=x+y\) has \(bc^{(1)}\in
F_{k-n-2}\). If \(z=bc^{(r)}\in F_p\), \(p=k-n-r-1\), then \(z\) is a cycle
and \(p\ne(k-1)-n\), so (P) gives \(x',y'\) of order \(\le p+1<k-n\) with
\(z-b(x'+y')\in F_{p-1}\). Put \(c^{(r+1)}=c^{(r)}-x'-y'\). By completeness
\(c=\lim c^{(r)}\) exists, \(bc=0\) and \(\Theta[c]=[x]\).

\emph{Degeneration.} \(d^r\) maps \(E^r_{p,k}\) to \(E^r_{p-r,k-1}\). Both
can be non-zero only if \(p=k-n\) and \(p-r=k-1-n\), that is, \(r=1\).
\end{proof}

This is Proposition~\ref{prop:convergence} with its proof completed.
Parts (1) and (2) need no completeness and hold in direct limits. The
chain model of Section~\ref{subsec:chains} is \emph{not} complete. In
\(C_1(\fA)=\varinjlim_N\fA_N^{\hot2}\) the series
\(\sum_{j\ge1}r^j\otimes r^{-2j}\) has total orders \(-j\to-\infty\), but
its first entries have unbounded order, so it lies in no
\(\fA_N^{\hot2}\). Following \cite[\S``Convergence and
degeneration'']{HHb}, we restore completeness with mixed complexes.

\subsubsection*{Mixed complexes}
Let \(\fA\) be one of
\begin{itemize}
\item the symbol algebra \(\cA_M\) of a closed manifold \(M\)
(Section~\ref{sec:closed});
\item \(\As\), or \(\Ads\);
\end{itemize}
with \(n\) the dimension of the underlying manifold. Let \(\mathfrak
B\subset\fA\) be the subalgebra of symbols of order \(\le0\) without
poles (for \(\Ads\), the Taylor image of that subalgebra of \(\As\)). Let
\(\fA_{L,M}\) be the elements with pole order \(\le L\) and order \(\le
M\) (\(L=0\) for \(\cA_M\)). The \emph{mixed chains} are
\[
 C_k(\mathfrak B,\fA)=\varinjlim_{L,M}S_{L,M,k},\qquad
 S_{L,M,k}=\fA_{L,M}\hot\mathfrak B^{\hot k},
\]
with the Hochschild differential with coefficients in the
\(\mathfrak B\)-bimodule \(\fA\): \(d_0\) and \(d_k\) put \(a_0b_1\) and
\(b_ka_0\) in slot \(0\), and the inner faces multiply in \(\mathfrak B\).

\begin{lemma}[Adapted from {\cite[Lemma ``A mixed complex for Laurent
complete symbols'']{HHb}}]\label{ha:lem:mixed}
\textup{(a)} Each \(S_{L,M}\) is a subcomplex. With the order filtration
it is complete in the sense of Proposition~\ref{ha:prop:algconv}.
\textup{(b)} The inclusion \(C(\mathfrak B,\fA)\to C(\fA)\) is a
quasi-isomorphism.
\end{lemma}

\begin{proof}
(a) Let \(T\in\mathfrak B\) be a symbol of order \(-1\) without poles,
elliptic with principal symbol \(q^{-1}\) (from (B3)). Its inverse in
\(\fA\) has order \(1\) and no poles. Put \(u=u_{L,M}=x^{-L}T^{-M}\). By
(B1), conjugation by \(x^{\pm L}\) preserves orders and smooth
coefficients. (In the cusp calculus \([x,P]\in x^2\Psic^{m-1}\), even
better than in the b-calculus.) Hence \(\fA_{L,M}=\mathfrak Bu=u\mathfrak
B\): for \(a\in\fA_{L,M}\), \(au^{-1}=x^{-L}(x^La)T^Mx^L\) is the conjugate
by \(x^{-L}\) of a pole-free symbol of order \(\le0\). So \(\fA_{L,M}\) is
a \(\mathfrak B\)-bimodule and \(S_{L,M}\) is a subcomplex. A quantization
(B1) identifies \(\fA_{L,M}\) and \(\mathfrak B\) with products of their
homogeneous components, and the product with the star product
\(\sum_lP_l\), \(P_l\) lowering the order by \(l\). By (F3), \(S_{L,M,k}\)
is the product over multidegrees \((j_0,\dots,j_k)\), \(j_0\le M\),
\(j_i\le0\), of the homogeneous tensor components. At a fixed total order
only finitely many multidegrees occur. So \(S_{L,M}\) consists of all
series in the total order bounded above by \(M\), and \(b\) commutes with
convergent sums.

(b) The argument is that of \cite{HHb}, which uses nothing specific to
the b-calculus; we give it in full. Write \(\fA_N=\fA_{N,N}\) and
\(u_N=x^{-N}T^{-N}\) (\(u_N=T^{-N}\) for \(\cA_M\)). By (a),
\(\fA_N=\mathfrak Bu_N\), and \(u_N^{-1}=T^Nx^N\in\mathfrak B\). The
chain spaces are \(C_k(\fA)=\varinjlim_N\fA_N^{\hot(k+1)}\) and
\(C_k(\mathfrak B,\fA)=\varinjlim_N\fA_N\hot\mathfrak B^{\hot k}\).

\emph{Complements.} \(\fA_{N-1}\) has a closed complement \(V_N\) in
\(\fA_N\) (\(V_0=\mathfrak B\)). For \(\Ads\), whose stages are products
of their homogeneous components \(x^jq^m\CI(Z)\), take the product of
the components with \(m=N\) or \(j=-N\). For \(\As\), take the
components of order \(N\) together with
\(x^{-N}\chi_0\pi_0^*(f|_{x=0})\), \(f\) running over the coefficients of
order \(<N\), with \(\pi_0\) the collar projection and \(\chi_0\) a
cut-off; for \(\cA_M\), the components of order \(N\). Then
\(\fA_N^{\hot k}=\bigoplus_{N_i\le N}V_{N_1}\hot\cdots\hot V_{N_k}\), a
finite direct sum of closed subspaces, and \(W_{k,N}=\bigoplus_{\max
N_i=N}V_{N_1}\hot\cdots\hot V_{N_k}\) gives
\(\varinjlim_N\fA_N^{\hot k}=\bigoplus_NW_{k,N}\) algebraically.

\emph{Two resolutions.} Let \(Q_k=\varinjlim_N\fA_N^{\hot(k+2)}\) and
\(P_k=\varinjlim_N\fA_N\hot\mathfrak B^{\hot k}\hot\fA_N\subset Q_k\)
(\(k\ge0\)), with the bar differential \(b'\), and put
\(P_{-1}=Q_{-1}=\fA\) with \(b'=\mu\) the multiplication \(Q_0\to\fA\).
These are complexes of \(\fA\)-bimodules (acting on the outer slots), and
\(b'\) is a bimodule map; \(P\) is a subcomplex because \(\mathfrak B\) is
a subalgebra. The left contraction \(s(a_0\otimes\dots\otimes a_{k+1})=
1\otimes a_0\otimes\dots\otimes a_{k+1}\), \(s(a)=1\otimes a\) on
\(Q_{-1}\), satisfies \(b's+sb'=1\) on \(Q\). For \(P\) let
\(P^{(N)}\subset P\) be the subcomplex of chains whose last slot lies in
\(\fA_N\) (\(b'\) maps that slot into \(\mathfrak B\fA_N=\fA_N\)), with
\(P^{(N)}_{-1}=\fA\). On \(P^{(N)}\) put
\[
 h_N(a_0\otimes\beta\otimes a)=(-1)^{k+1}a_0\otimes\beta\otimes
 au_N^{-1}\otimes u_N,\qquad h_N(a)=au_N^{-1}\otimes u_N\ \ (a\in
 P_{-1}),
\]
for \(\beta\in\mathfrak B^{\hot k}\), \(a\in\fA_N\). Here
\(au_N^{-1}\in\mathfrak B\), so \(h_N\) maps \(P^{(N)}_k\) to
\(P^{(N)}_{k+1}\). A direct computation gives \(b'h_N+h_Nb'=1\) on
\(P^{(N)}\): the faces of \(b'h_N(a_0\otimes\beta\otimes a)\) that do not
involve the last two slots cancel those of \(h_Nb'\), the face merging
\(au_N^{-1}\) with \(u_N\) gives \(a_0\otimes\beta\otimes a\), and the
face merging \(\beta\)'s last entry \(b_k\) with \(au_N^{-1}\) cancels
\(h_N\) of the last face of \(b'\), because \((b_ka)u_N^{-1}=b_k(au_N^{-1})\).
(It is important that \(u_N\) is the same in both terms; this is why we
work in \(P^{(N)}\).)

\emph{Comparison maps.} Let \(\iota:P\to Q\) be the inclusion. We
construct a bimodule chain map \(g:Q\to P\) over the identity of \(\fA\).
Put \(g_{-1}=1\) and \(g_0=1\) (\(P_0=Q_0\)). Let \(k\ge1\), and suppose
\(g_{k-1}\) is constructed, continuous on stages, and maps the
generators \(1\otimes W_{k-1,N}\otimes1\) into \(P^{(\nu(k-1,N))}\) for
some \(\nu(k-1,N)\). For \(w\in W_{k,N}\), \(b'(1\otimes w\otimes1)\) is a
sum of terms \(a\cdot(1\otimes w'\otimes1)\cdot a'\) with \(a,a'\in\fA_N\)
and \(w'\in\fA_{2N}^{\hot(k-1)}\), so \(z(w)=g_{k-1}b'(1\otimes
w\otimes1)\) lies in \(P^{(\nu)}\) with
\(\nu=\nu(k,N)=N+\max_{N'\le2N}\nu(k-1,N')\). It is a cycle of the
augmented complex: \(b'z(w)=g_{k-2}b'b'(\cdot)=0\) for \(k\ge2\), and
\(\mu z(w)=\mu b'(1\otimes w\otimes1)=0\) for \(k=1\). Put
\(g_k(1\otimes w\otimes1)=h_\nu(z(w))\), so that
\(b'g_k(1\otimes w\otimes1)=z(w)\), and extend by
\(g_k(a\otimes w\otimes a')=a\,g_k(1\otimes w\otimes1)\,a'\). Since
multiplication is continuous on stages, this is continuous on each
\(\fA_N\hot W_{k,N}\hot\fA_N\), hence well defined on \(Q_k\). Both
\(b'g_k\) and \(g_{k-1}b'\) are bimodule maps that agree on generators,
so \(g\) is a chain map. In the same way one builds bimodule homotopies
\(K:Q_k\to Q_{k+1}\) with \(b'K+Kb'=\iota g-1\), by
\(K_k(1\otimes w\otimes1)=s\bigl((\iota g-1-K_{k-1}b')(1\otimes
w\otimes1)\bigr)\) (\(K_{-1}=K_0=0\)), and \(H:P_k\to P_{k+1}\) with
\(b'H+Hb'=g\iota-1\), by the same formula with \(h_\nu\) in place of
\(s\), \(\nu\) being chosen, as for \(g\), so that the argument lies in
\(P^{(\nu)}\); the generators of \(P_k\) are \(1\otimes\beta\otimes1\),
\(\beta\in\mathfrak B^{\hot k}=W_{k,0}\). In each case the argument of
\(s\) or \(h_\nu\) is a cycle by the inductive hypothesis.

\emph{Descent.} The map \(\rho(a_0\otimes w\otimes a)=aa_0\otimes w\)
maps \(Q_k\) onto \(C_k(\fA)\) and \(P_k\) onto \(C_k(\mathfrak B,\fA)\),
and \(\rho b'=b\rho\) (the face merging the last two slots of the bar
complex becomes the last face \(d_k\) of \(b\)). Moreover
\(\rho(aX)=\rho(Xa)\) for \(a\in\fA\). For a bimodule map \(F\) between
\(P\) and \(Q\) (in either direction) define \(\bar F(a_0\otimes w)=
\rho F(a_0\otimes w\otimes1)\). Then \(\bar F\rho=\rho F\): on elementary
tensors, \(\rho F(a_0\otimes w\otimes a)=\rho(F(a_0\otimes w\otimes1)a)
=\rho(aF(a_0\otimes w\otimes1))=\bar F(aa_0\otimes w)\), and both sides
are continuous. Hence \(\bar g\) is a chain map \(C(\fA)\to
C(\mathfrak B,\fA)\), \(\bar\iota\) is the inclusion, and \(\bar K\),
\(\bar H\) are homotopies \(\bar\iota\bar g\simeq1\),
\(\bar g\bar\iota\simeq1\). So the inclusion is a chain homotopy
equivalence.
\end{proof}

\begin{lemma}[Fixed-stage \(E^1\)]\label{ha:lem:stageE1}
Let \(q\le M-1\). The \(b_0\)-homology of \(S_{L,M}\) at total order \(q\)
is, via \(\chi\), the space of homogeneous \(k\)-forms of degree \(q\) on
the cotangent bundle minus its zero section, with coefficients
\begin{itemize}
\item smooth, for \(\cA_M\);
\item in \(x^{-L}\Om\), i.e.\ \(x^{-L}\) times \emph{ordinary} smooth forms
on \(\cT\setminus0\) (generated by \(dx,dy,d\tau,d\eta\)), for \(\As\);
\item in the formal analogue, for \(\Ads\).
\end{itemize}
The inclusions of stages induce the obvious inclusions of forms. In the
direct limit one gets all smooth, Laurent and formal Laurent homogeneous
forms, which is the \(E^1\) of Proposition~\ref{prop:As}.
\end{lemma}

\begin{proof}
The associated graded algebra of \(\mathfrak B\) is \(\CI(Y)[t]\), with
\(Y\) the cosphere bundle, \(t=r^{-1}\) of weight \(-1\), and \(\CI(Y)\)
replaced by its formal version for \(\Ads\). The graded coefficient module
is \(x^{-L}t^{-M}\CI(Y)[t]\), free of rank one with equal left and right
actions. So its Hochschild complex is that of \(\CI(Y)[t]\) shifted by
\(x^{-L}t^{-M}\), and Corollary~\ref{ha:cor:homog} gives \(x^{-L}t^{-M}\)
times forms in \(\Om(Y)[t]\oplus\Om(Y)[t]\,dt\). A homogeneous form of
degree \(q\) is \(t^{-q}(\alpha-\beta\wedge dt/t)\), and lies there if and
only if \(M-q\ge0\) and \(M-q-1\ge0\). The forms on \(Y\) are ordinary
forms, since \(\mathfrak B\) has smooth coefficients.
\end{proof}

\subsubsection*{The stage argument}

\begin{proposition}[Adapted from {\cite[\S``Convergence and
degeneration'']{HHb}}]\label{ha:prop:stage}
Let \(D\subset C(\mathfrak B,\fA)\) be a subcomplex with closed
subcomplexes \(D_{L,M}=D\cap S_{L,M}\), and suppose:
\begin{enumerate}
\item[(E2)] the \(E^2\) term of \(D\) vanishes off the line \(p=k-n\);
\item[(St)] for \(L\ge L_0\) and every \(M\), step (P) of
Proposition~\ref{ha:prop:algconv} can be carried out inside \(D_{L,M}\)
whenever \(p\le M-2\) and \(p\ne k-n\).
\end{enumerate}
Then \(\Theta:H_k(D)\to E^2_{k-n,k}(D)\) is an isomorphism, and
\(d^r=0\) for \(r\ge2\).
\end{proposition}

\begin{proof}
Parts (1) and (2) of Proposition~\ref{ha:prop:algconv} hold in the direct
limit. For injectivity, let \(c\in F_{k-n}\) be a cycle with
\([c_{k-n}]=0\). Perform one step (P) in the direct limit and choose a
stage \(D_{L,M}\) containing the remainder, with \(L\ge L_0\) and \(M\)
so large that all later orders are \(\le M-2\). All later steps are off
the critical line, so by (St) they stay in \(D_{L,M}\), and the infinite
sum converges there by Lemma~\ref{ha:lem:mixed}(a). For surjectivity,
choose \(x\) and \(y\) in the direct limit, then a stage containing them,
and continue in the same way. Degeneration is as before.
\end{proof}

\begin{theorem}[Convergence for \(\cA_M\) and \(\Is\)]\label{ha:thm:conv}
For the symbol algebra \(\cA_M\) of a closed manifold, and for \(\Is\),
the order spectral sequence converges and degenerates at \(E^2\).
\(\Theta\) is an isomorphism \(\HH_k\to E^2_{k-n,k}\), given by the leading
symbol of a representative in \(F_{k-n}\).
\end{theorem}

\begin{proof}
\emph{\(\cA_M\).} Take \(D=C(\mathfrak B,\cA_M)\), which computes
\(\HH(\cA_M)\) by Lemma~\ref{ha:lem:mixed}(b). (E2) is
Propositions~\ref{prop:d1}, \ref{prop:star} and~\ref{prop:homog}. For
(St), let \(\theta\) be the HKR form of \(c_p\). It is homogeneous of
degree \(p\) and \(\delta\theta=0\). Then \(*\theta\) is closed and
homogeneous of degree \(w=p+n-k\ne0\). By Cartan's formula \(\eta=
\iota_R(*\theta)/w\) satisfies \(d\eta=*\theta\), and \(\zeta=*^{-1}\eta\)
is homogeneous of degree \(p+1\le M-1\) with \(\delta\zeta=\pm\theta\). By
Lemma~\ref{ha:lem:stageE1}, \(\pm i\zeta\) is the class of a \(b_0\)-cycle
\(x\) in the same stage at order \(p+1\), with \(d^1[x]=[c_p]\) (the
proof of Proposition~\ref{prop:d1} applies to mixed chains: it only uses
the leading symbols in each slot). Then \(c_p-b_1x\) has zero class in the
stage \(E^1\) at order \(p\le M-1\), so it is \(b_0y\) with \(y\) in the
stage.

\emph{\(\Is\).} \(\Is\) is H-unital (Proposition~\ref{prop:Is}). By
Theorem~\ref{ha:thm:excision}, \(C(\Is)\) is quasi-isomorphic to the kernel
\(\cK\) of the Taylor map \(C(\As)\to C(\Ads)\). Let \(D\) be the kernel
of the Taylor map \(C(\mathfrak B,\As)\to C(\mathfrak B_\partial,\Ads)\).
The Taylor maps are onto on every stage, by (F1) and Borel's lemma. So
Lemma~\ref{ha:lem:mixed}(b), for \(\As\) and \(\Ads\), and the five lemma
show that \(D\to\cK\) is a quasi-isomorphism. The Taylor map is also onto
on the stage \(E^1\) (Borel's lemma for forms). So the stage \(E^1\) of
\(D\) at orders \(q\le M-1\) is the space of all flat homogeneous forms
of degree \(q\). The operators \(*\), \(*^{-1}\) and \(\iota_R\) preserve
flat forms: \(*\) only introduces the factors \(x^{\pm2}\). So the
argument for \(\cA_M\) applies without change and gives (St). (E2) is the
\(E^2\) computation of Proposition~\ref{prop:Is}. The maps between the
three complexes preserve the order and induce the flat HKR
identification on graded pieces. So the resulting isomorphism is the
leading-symbol map.
\end{proof}

\subsubsection*{The Laurent algebras \(\As\), \(\Ads\) and the cusp twist}
In \cite{HHb} condition (St) is verified for the b-calculus analogues of
\(\As\) and \(\Ads\) \cite[Lemma ``Lifts in a fixed coefficient
stage'']{HHb}. The cusp version fails. (Convergence nevertheless holds;
it is proved by a different argument in Appendix~\ref{su:sec}.) By Lemma~\ref{ha:lem:stageE1} the
stage \(E^1\) of \(\As\) consists of \(x^{-L}\) times ordinary forms. The
symplectic star, however, preserves \(x^{-L}\) times \emph{cusp} forms.

\begin{example}[The stage obstruction, \(n=1\)]\label{ha:ex:obstruction}
Let \(n=1\) and work near one boundary point, with coordinates \((x,\tau)\)
and \(\{a,b\}=\partial_\tau a\,x^2\partial_xb-x^2\partial_xa\,
\partial_\tau b\). The critical line is \(p=k-1\). Fix \(L\), \(M\) and
\(p\le M-2\) with \(p\ne-1\). The \(0\)-form \(\theta=x^{-L}|\tau|^p\) lies
in the stage \(E^1_{p,0}\), and \(d^1\theta=0\). The stage \(1\)-forms of
degree \(p+1\) are \(\zeta=x^{-L}(a\,d\tau+c\,dx)\) with \(a,c\) smooth and
homogeneous of degrees \(p\) and \(p+1\). For \(1\)-forms
\(\delta(f\,dg)=\{f,g\}\), so
\[
 \delta\zeta=\{x^{-L}a,\tau\}+\{x^{-L}c,x\}=-x^2\partial_x(x^{-L}a)+
 x^{2-L}\partial_\tau c\ \in\ x^{1-L}\CI .
\]
Hence \(\theta\) is not in the image, and \([\theta]\ne0\) in the stage
\(E^2_{p,0}\), although \([\theta]=0\) in the direct limit:
\(\theta=\delta\bigl(x^{-L-1}|\tau|^p\,d\tau/(L+1)\bigr)\). So (St) fails in
every stage. In the b-calculus the same computation, with
\(x\partial_\xi\wedge\partial_x\) in place of \(x^2\partial_\tau\wedge
\partial_x\), gives \(-x\partial_x(x^{-L}a)=Lx^{-L}a+O(x^{1-L})\). This
reaches the leading coefficient when \(L\ne0\). That is the invertibility
of \(\mathcal L_{x\partial_x}=-L\) on the leading coefficient used in
\cite{HHb}. In the cusp case its replacement \(\mathcal L_{x^2\partial_x}\)
raises the power of \(x\), so it is nilpotent on the leading coefficient.

This obstructs the \emph{method} (St), not convergence itself. For
every \(n\), Theorem~\ref{su:thm:B10} proves that \(\Theta\) is an
isomorphism, and for \(n=1\) so does Proposition~\ref{ha:prop:xadic}
below. There the algebra is graded by the power of \(x\) plus the
order, each correction step leaves a remainder whose lowest power of
\(x\) is strictly higher than before, and so a cycle that starts two pole
levels inside a stage never produces a term at the lowest power
\(x^{-L}\), which is where (St) fails.
\end{example}

The same computation in general dimension shows that \(\delta\) of a stage
form never reaches the leading coefficient \(x^{-L}\) through the
\((x,\tau)\)-part of the Poisson structure. On leading coefficients only
the boundary part \(\partial_\eta\wedge\partial_y\) acts, and for it
\(\tau\) is a Casimir. So the stage \(E^2\) off the critical line contains
classes such as \(x^{-L}|\tau|^p\) in every stage. We do not need a
precise description. Under the transcendental change \(\rho=e^{-1/x}\), which
turns the cusp structure into a b-structure
(Section~\ref{sec:cusp}), one has \(x^{-L}=(\log1/\rho)^L\). So the cusp
Laurent algebras are b-algebras with \emph{logarithmic} coefficients, on
which \(\rho\partial_\rho\) acts nilpotently. What survives is the
following.

\begin{lemma}\label{ha:lem:WL}
Let \(\theta\in x^{-L}\Om^k\) be homogeneous of degree \(p\ne k-n\) with
\(\delta\theta=0\). Then \(\theta=\delta\zeta\) with \(\zeta\in
x^{-L-2}\Om^{k+1}\) homogeneous of degree \(p+1\).
\end{lemma}

\begin{proof}
Let \(\Om_c\) be the smooth cusp forms, generated by \(dx/x^2,dy,d\tau,
d\eta\). Then \(\Om\subset\Om_c\subset x^{-2}\Om\). The cusp symplectic form
is a smooth non-degenerate cusp \(2\)-form, so \(*\) and \(*^{-1}\) preserve
\(x^{-L}\Om_c\), and so does \(\iota_R\). The construction in the proof of
Theorem~\ref{ha:thm:conv} gives \(\zeta\in x^{-L}\Om_c\subset
x^{-L-2}\Om\).
\end{proof}

\begin{proposition}[Partial convergence for \(\As\) and \(\Ads\)]
\label{ha:prop:laurent}
Let \(\fA=\As\) or \(\Ads\), with the mixed complex
\(C(\mathfrak B,\fA)\), which computes \(\HH(\fA)\).
\begin{enumerate}
\item Every class in \(\HH_k(\fA)\) has a representative in
\(F_{k-n}\), and \(\Theta:\HH_k(\fA)\to E^2_{k-n,k}\) is well defined.
Here \(E^2\) is the one computed in Propositions~\ref{prop:As}
and~\ref{prop:Ads}.
\item \(\ker\Theta=\bigcap_pF_p\HH_k(\fA)\).
\item Let \(x\in S_{L,M}\) be a \(b_0\)-cycle of order \(k-n\le M-2\)
with \(d^1[x]=0\). For every \(r\ge1\) there is \(c^{(r)}\in
S_{L+2r,M}\) with leading term \(x\) and \(bc^{(r)}\in F_{k-n-r-1}\).
\item \(\Theta\) is an isomorphism if and only if \(\HH_k(\fA)\to
H_k(\widehat C)\) is one. Here \(\widehat C\) is the complex of series
\(\sum_{p\le P}c_p\) in which each \(c_p\) lies in some stage, with no
uniform bound on the poles.
\end{enumerate}
\end{proposition}

\begin{proof}
(1) and (2) are Proposition~\ref{ha:prop:algconv}(1),(2), which need no
completeness. For (3), \([b_1x]=0\) already in the stage \(E^1\), by the
injectivity in Lemma~\ref{ha:lem:stageE1}. Each later step (P) is off the
critical line. By Lemma~\ref{ha:lem:WL} its primitive lies in
\(x^{-L'-2}\Om\), so the corrections lie in the stage \(L'+2\). For (4):
\(\widehat C\) is a complex because each \(b_l\) preserves the stages. It
is complete, and its graded pieces and \(E^2\) agree with those of
\(C(\mathfrak B,\fA)\). So \(H_k(\widehat C)\cong E^2_{k-n,k}\) by
Proposition~\ref{ha:prop:algconv}, compatibly with \(\Theta\).
\end{proof}


\begin{remark}[The last step of the convergence, (B10)]\label{ha:rem:quoted}
Proposition~\ref{ha:prop:laurent} reduces the convergence for \(\As\) and
\(\Ads\) to one statement. Corrections whose poles grow by up to two per
order (part (3)) must be replaceable by corrections with a uniform pole
bound. Equivalently, there are no phantom classes
\(\bigcap_pF_p\HH_k\ne0\), and the towers in (3) can be completed in
\(C(\mathfrak B,\fA)\). By Proposition~\ref{ha:prop:equiv} below it is
enough to prove this for one of the two algebras. MN assert it (``as in
the boundaryless case'', \cite[p.~12]{MN}). It is item (B10) of
Section~\ref{subsec:background}, and it is now proved for every \(n\):
for \(n=1\) by Proposition~\ref{ha:prop:xadic} and
Corollary~\ref{ha:cor:n1}, and for all \(n\) by Theorem~\ref{su:thm:B10},
which follows Lauter and Moroianu \cite{LauterMoroianu}. Three other
routes do not work: the fixed-stage argument of \cite{HHb}
(Example~\ref{ha:ex:obstruction}), the \(x\)-adic filtration as a whole,
whose \(E^1\) term needs \(\HH\) of the suspended symbol algebra
(Remark~\ref{ha:rem:xadicfail}), and an import from Benameur--Nistor
\cite{BN} (Remark~\ref{ha:rem:BN}). The proof of
Theorem~\ref{su:thm:B10} uses the \(x\)-grading only for the part of a
remainder with one extra pole, and it needs from the suspended symbol
algebra only the degeneration of its order spectral sequence, not its
homology.
\end{remark}

\subsection{The boundary connecting map (B9)}\label{ha:subsec:B9}

MN compute the connecting map of \(0\to\Is\to\As\to\Ads\to0\) with a
multiplicative splitting \(\Ads\to\As\), which we could not construct
(Remark~\ref{rem:s0}). The leading part of a connecting map only needs a
\emph{linear} lift, chosen so that its top component is a \(b_0\)-cycle.
The following lemma is stated for complexes of the kind considered in
Proposition~\ref{ha:prop:algconv}; it uses no completeness.

\begin{lemma}[Leading part of a connecting map]\label{ha:lem:connlead}
Let \(0\to K\xrightarrow{\iota}C\xrightarrow{\pi}C''\to0\) be chain maps
between complexes of series as in Proposition~\ref{ha:prop:algconv}, with
\(b=\sum_{l\ge0}b_l\) in each. Suppose that \(\iota\) and \(\pi\) map
\(C_{p,k}\)-components to \(C_{p,k}\)-components, that
\(0\to K_{p,k}\to C_{p,k}\to C''_{p,k}\to0\) is exact for all \(p,k\), that
\(\pi\) maps \(F_pC_k\) onto \(F_pC''_k\) for every \(p\), and that
\(\pi_*:E^1_{p,k}(C)\to E^1_{p,k}(C'')\) is onto for all \(p,k\).
\begin{enumerate}
\item The sequence \(0\to E^1(K)\to E^1(C)\to E^1(C'')\to0\) is exact.
It is a short exact sequence of complexes for \(d^1\); let
\(\partial^1:E^2_{p,k}(C'')\to E^2_{p-1,k-1}(K)\) be its connecting map.
\item Let \(c\in F_pC''_k\) be a cycle. There is \(\tilde c\in F_pC_k\) with
\(\pi\tilde c=c\) and \(b_0\tilde c_p=0\). For every such \(\tilde c\),
the chain \(b\tilde c\) is a cycle of \(K\) lying in \(F_{p-1}K_{k-1}\),
its component \((b\tilde c)_{p-1}\) is a \(b_0\)-cycle, and its class in
\(E^2_{p-1,k-1}(K)\) is \(\partial^1\) of the class of \(c_p\) in
\(E^2_{p,k}(C'')\).
\item Suppose moreover that \(E^2(K)\) and \(E^2(C'')\) vanish off the
line \(p=k-n\), and let \(\Theta_K\), \(\Theta_{C''}\) be the maps of
Proposition~\ref{ha:prop:algconv}(1). Then, with \(\partial\) the
connecting map of the long exact homology sequence,
\[
 \Theta_K\circ\partial=\partial^1\circ\Theta_{C''}:\ H_k(C'')\to
 E^2_{k-1-n,k-1}(K).
\]
\end{enumerate}
\end{lemma}

\begin{proof}
(1) For fixed \(p\) the long exact sequence of
\(0\to(K_{p,\bullet},b_0)\to(C_{p,\bullet},b_0)\to(C''_{p,\bullet},b_0)\to0\)
has zero connecting maps, because \(\pi_*\) is onto in every degree. The
maps commute with \(d^1\), since \(\iota\) and \(\pi\) commute with \(b\)
and preserve components, hence commute with each \(b_l\).

(2) Since \(c\) is a cycle, \(b_0c_p=0\). Choose a \(b_0\)-cycle \(z\in
C_{p,k}\) with \(\pi_*[z]=[c_p]\) in \(E^1(C'')\). Then \(c_p-\pi z=b_0w''\)
with \(w''\in C''_{p,k+1}\). Lift \(w''\) to \(w\in C_{p,k+1}\) and put
\(\tilde c_p=z+b_0w\). Then \(b_0\tilde c_p=0\) and \(\pi\tilde c_p=c_p\).
Lift \(c-c_p\in F_{p-1}C''_k\) to \(t\in F_{p-1}C_k\), and put
\(\tilde c=\tilde c_p+t\).

Now let \(\tilde c\) be any lift with \(b_0\tilde c_p=0\). Since
\(\pi b\tilde c=bc=0\), the chain \(b\tilde c\) lies in \(K\), and it is a
cycle. Its component of order \(p\) is \(b_0\tilde c_p=0\), so \(b\tilde
c\in F_{p-1}\), and its component of order \(p-1\) is
\(e=b_1\tilde c_p+b_0\tilde c_{p-1}\). The order-\((p-1)\) component of
\(b(b\tilde c)=0\) is \(b_0e\), so \(e\) is a \(b_0\)-cycle. In
\(E^1(C)\), \(\iota_*[e]=[b_1\tilde c_p]=d^1[\tilde c_p]\). Here
\([\tilde c_p]\) lifts \([c_p]\), and \([c_p]\) is a \(d^1\)-cycle because
\(b_1c_p=-b_0c_{p-1}\). By the definition of the connecting map (the snake
lemma), \([e]\) represents \(\partial^1\) of the class of \(c_p\).

(3) By Proposition~\ref{ha:prop:algconv}(1) a class of \(H_k(C'')\) has a
representative \(c\in F_{k-n}\). By the snake lemma \(\partial[c]=[b\tilde
c]\) for any lift \(\tilde c\). Take \(\tilde c\) as in (2) with \(p=k-n\).
Then \(b\tilde c\in F_{(k-1)-n}K_{k-1}\) represents \(\partial[c]\), and
its leading component represents \(\partial^1[c_{k-n}]\) by (2).
\end{proof}

Nothing in the lemma is multiplicative. The lift is constructed through
\(E^1\), by Borel's lemma in the application. The fact that the lift is
not an algebra map affects only the components of \(b\tilde c\) of
order \(\le p-2\), and those do not enter \(\Theta_K\).

\begin{theorem}[{The boundary connecting map; \cite[Prop.~5(3)]{MN}}]
\label{ha:thm:bdry}
Let \(\partial:\HH_k(\Ads)\to\HH_{k-1}(\Is)\) be the connecting map of
\(0\to\Is\to\As\to\Ads\to0\). Let \(\Theta_{\partial,\sigma}:\HH_k(\Ads)\to
E^2_{k-n,k}(\Ads)\) be the leading-symbol map of
Proposition~\ref{ha:prop:laurent}(1), and \(\Theta_{\Is}:
\HH_{k-1}(\Is)\to E^2_{k-1-n,k-1}(\Is)\) the isomorphism of
Theorem~\ref{ha:thm:conv}. Identify, through the symplectic star,
\[
 E^2_{k-n,k}(\Ads)\cong H^{2n-k}(Z\times\Ssig)\oplus H^{2n-k-1}(Z\times\Ssig)
 \,d\theta_\partial,\qquad
 E^2_{k-1-n,k-1}(\Is)\cong\Hrel^{2n-k+1}(Y\times\Ssig),
\]
as in Propositions~\ref{prop:Ads} and~\ref{prop:Is}. Then
\[
 \Theta_{\Is}\circ\partial=(-1)^k\,i\,\bigl(\delta_{\mathrm{dR}}\oplus0
 \bigr)\circ\Theta_{\partial,\sigma},
\]
where \(\delta_{\mathrm{dR}}:H^{2n-k}(Z\times\Ssig)\to\Hrel^{2n-k+1}(Y
\times\Ssig)\) is the connecting map of the pair. No convergence
statement for \(\As\) or \(\Ads\) is used.
\end{theorem}

\begin{proof}
\emph{The complexes.} Let \(C=C(\mathfrak B,\As)\) and \(C''=C(\mathfrak
B_\partial,\Ads)\) be the mixed complexes, let \(\pi:C\to C''\) be the
Taylor map, and let \(D=\ker\pi\), as in the proof of
Theorem~\ref{ha:thm:conv}. Choose the quantization of \(\Ads\) induced by
that of \(\As\). Then \(\pi\) takes the homogeneous component of order
\(p\) of a chain to the homogeneous component of order \(p\) of its image,
and \(D_{p,k}\) is the kernel on \(C_{p,k}\). The map \(\pi\) is onto on
every stage, by (F1) and Borel's lemma, hence onto \(F_pC''_k\) and onto
each \(C''_{p,k}\). The inclusions of \(0\to D\to C\to C''\to0\) into
\(0\to\cK\to C(\As)\to C(\Ads)\to0\) are quasi-isomorphisms
(Lemma~\ref{ha:lem:mixed}(b) and the five lemma), and \(C(\Is)\to\cK\) is
one by Theorem~\ref{ha:thm:excision}. By naturality of the long exact
sequences, \(\partial\) is the connecting map of \(0\to D\to C\to
C''\to0\), and \(\Theta_{\Is}\) is computed in \(D\)
(Theorem~\ref{ha:thm:conv}).

\emph{The \(E^1\) and \(E^2\) terms.} By Lemma~\ref{ha:lem:stageE1} and
the proof of Theorem~\ref{ha:thm:conv}, in the direct limit over stages
\(\chi\) identifies \(E^1_{p,k}\) of \(D\), \(C\), \(C''\) with the flat,
Laurent and formal Laurent homogeneous \(k\)-forms of degree \(p\) on
\(\cT\setminus0\), and \(\pi_*\) with Taylor expansion at \(x=0\). Taylor
expansion of forms is onto by Borel's lemma. So
Lemma~\ref{ha:lem:connlead} applies. By Proposition~\ref{prop:d1}
(applied to mixed chains as in the proof of Theorem~\ref{ha:thm:conv}),
\(d^1=-i\delta\) in all three complexes. The connecting map of
\((-i\delta)\)-complexes is \(-i\) times that of \(\delta\)-complexes, so
\[
 \partial^1[\mu]=-i\,\partial_\delta[\mu],\qquad
 \partial_\delta[\mu]=[\delta\lambda],
\]
for any Laurent homogeneous form \(\lambda\) whose Taylor series is
\(\mu\). The \(E^2\) terms of \(D\) and \(C''\) vanish off the line
\(p=k-n\) (Propositions~\ref{prop:Is} and~\ref{prop:Ads}), so
Lemma~\ref{ha:lem:connlead}(3) gives \(\Theta_{\Is}\circ\partial=
-i\,\partial_\delta\circ\Theta_{\partial,\sigma}\).

\emph{From \(\delta\) to \(d\).} The symplectic star is defined by the
cusp symplectic form, which is smooth and non-degenerate as a cusp form.
So it acts on flat, Laurent and formal Laurent forms, and it commutes with
Taylor expansion. On \(k\)-forms \(\delta=(-1)^{k-1}*^{-1}d\,*\)
(Proposition~\ref{prop:star}). If \(\mu\) is a formal Laurent \(k\)-form
with \(\delta\mu=0\) and \(\lambda\) a Laurent lift, then \(*\lambda\) is a
Laurent lift of \(*\mu\), and \(*\delta\lambda=(-1)^{k-1}d(*\lambda)\).
Hence \(*\,\partial_\delta[\mu]=(-1)^{k-1}\partial_d[*\mu]\), where
\(\partial_d\) is the connecting map of the de Rham complexes of
homogeneous forms of degree \(0\) (the degree of \(*\mu\) for \(\mu\) on
the critical line). Altogether, with the identifications through \(*\),
\(\Theta_{\Is}\circ\partial=-i(-1)^{k-1}\partial_d\circ
\Theta_{\partial,\sigma}=(-1)^ki\,\partial_d\circ\Theta_{\partial,\sigma}\).

\emph{The de Rham connecting map.} Homogeneous forms of degree \(0\) are
pulled back from \(Y\times\Ssig\) (Proposition~\ref{prop:homog}), and the
flat, Laurent and formal versions correspond to the three complexes of
Proposition~\ref{prop:dR} for \(M=Y\times\Ssig\), with boundary
\(Z\times\Ssig\); this is \eqref{eq:bH}. Let \(\pi_0\) be the projection
of the collar of \(Y\times\Ssig\) onto \(Z\times\Ssig\), and \(\chi_0\) a
cut-off equal to \(1\) near the boundary, as in the proof of
Proposition~\ref{prop:dR}. A class without \(d\theta_\partial=[dx/x]\) is
represented by the formal form \(\pi_0^*\beta\), with \(\beta\) closed on
\(Z\times\Ssig\). It has the Laurent lift \(\chi_0\pi_0^*\beta\), and
\(d(\chi_0\pi_0^*\beta)=d\chi_0\wedge\pi_0^*\beta\) is supported in the
interior. Its class in \(\Hrel\) is \(\delta_{\mathrm{dR}}[\beta]\), by
the definition of the connecting map of the pair (the flat complex
computes \(\Hrel\), Proposition~\ref{prop:dR}). A class
\(\beta\wedge dx/x\) has the Laurent lift \(\chi_0\pi_0^*\beta\wedge
dx/x\), which is closed because \(d\chi_0\) is a multiple of \(dx\) in
the collar. So \(\partial_d\) vanishes on it.
\end{proof}

\begin{proposition}[The two convergence statements are equivalent]
\label{ha:prop:equiv}
For every \(n\ge1\) the following are equivalent:
\begin{enumerate}
\item \(\Theta:\HH_k(\As)\to E^2_{k-n,k}(\As)\) is an isomorphism for all
\(k\);
\item \(\Theta:\HH_k(\Ads)\to E^2_{k-n,k}(\Ads)\) is an isomorphism for
all \(k\).
\end{enumerate}
If they hold, then Proposition~\ref{prop:boundaryLES} holds as stated,
with the constant \((-1)^ki\) in part (3).
\end{proposition}

\begin{proof}
Consider the ladder whose top row is the long exact sequence
\[
 \cdots\to\HH_k(\Is)\to\HH_k(\As)\to\HH_k(\Ads)\xrightarrow{\partial}
 \HH_{k-1}(\Is)\to\cdots
\]
(Theorem~\ref{ha:thm:excision}; \(\Is\) is H-unital by
Proposition~\ref{prop:Is}). The bottom row is the long exact sequence of
\(d^1\)-homology of the short exact sequence of complexes
\(0\to E^1(D)\to E^1(C)\to E^1(C'')\to0\) of
Lemma~\ref{ha:lem:connlead}(1). Since \(d^1\) maps \((p,k)\) to
\((p-1,k-1)\), this sequence splits into the sequences of the lines
\(p-k=\mathrm{const}\). We keep the line \(p=k-n\):
\[
 \cdots\to E^2_{k-n,k}(\Is)\to E^2_{k-n,k}(\As)\to E^2_{k-n,k}(\Ads)
 \xrightarrow{\partial^1}E^2_{k-1-n,k-1}(\Is)\to\cdots .
\]
It is exact. The vertical maps are the maps \(\Theta\). They are defined
by Proposition~\ref{ha:prop:laurent}(1) for \(\As\) and \(\Ads\), and
Theorem~\ref{ha:thm:conv} shows that \(\Theta\) is an isomorphism for
\(\Is\). The squares with the maps induced by \(\Is\to\As\to\Ads\)
commute: these maps preserve the order filtration and the components, so
they take a representative in \(F_{k-n}\) to one, and its leading term to
the leading term of the image. The square with the connecting maps
commutes by Lemma~\ref{ha:lem:connlead}(3), applied as in the proof of
Theorem~\ref{ha:thm:bdry}. The five lemma, applied with \(\HH_k(\As)\) or
with \(\HH_k(\Ads)\) in the middle, gives (2)\(\Rightarrow\)(1) and
(1)\(\Rightarrow\)(2). The last statement follows from the
naturality just used, for parts (1) and (2) of
Proposition~\ref{prop:boundaryLES}, and from Theorem~\ref{ha:thm:bdry}
for part (3).
\end{proof}

\subsection{Dimension one: the \texorpdfstring{\(x\)}{x}-adic filtration}
\label{ha:subsec:xadic}

For \(n=1\) the boundary symbol algebra is commutative, and the
\(x\)-adic filtration settles (B10). Let \(\partial X\) consist of \(N\)
points. Then \(Z\) consists of \(2N\) points \(z\), one for each boundary
point and each sign \(\varepsilon_z=\pm\) of \(\tau\). Put \(s=|\tau|^{-1}\).

\begin{lemma}[\(\Ads\) for \(n=1\)]\label{ha:lem:n1alg}
As a vector space, \(\Ads\) is the space of series \(\sum_{j\ge-L,\,r\ge
-M}c_{jr}x^js^r\) with \(c_{jr}\in\C^Z\). The product is determined by
\[
 (x^ja)(x^kb)=\sum_{\mu\ge0}c_{k,\mu}\,x^{j+k+\mu}(D^\mu a)\,b,\qquad
 a,b\in\C^Z\Lser{s},\qquad D=\varepsilon_z\,i\,s^2\partial_s\ \text{at }z,
\]
with constants \(c_{k,\mu}\) such that \(c_{k,0}=1\) and \(c_{k,1}=k\).
In particular \(\Ads\) is graded by \(g=J+p\), where \(J\) is the power
of \(x\) and \(p\) the order (\(x^js^r\) has \(J=j\), \(p=-r\)).
\end{lemma}

\begin{proof}
For a point, \(\Psis^{\Z}/\Psis^{-\infty}\) is, through \(\In\), the
algebra of classical symbols in \(\tau\) modulo \(\cS(\R)\), with the
pointwise product \eqref{eq:Inkernel}. Its elements are pairs of formal
Laurent series in \(s\), at \(\tau=+\infty\) and at \(\tau=-\infty\). Let
\(D=\mathrm{ad}_t\). By Lemma~\ref{lem:twist}(1), \(D\) acts on indicial
families by \(-i\partial_\tau\), which is \(\varepsilon_zis^2\partial_s\)
in the variable \(s\). By Lemma~\ref{lem:twist}(2),
\([x,a]=-x\,(Da)\,x\), that is, \(ax=xa+x(Da)x\). Iterating gives
\(ax=\sum_{\mu\ge0}x^{\mu+1}D^\mu a\), and by induction on \(k\ge0\) an
expansion \(ax^k=\sum_\mu c_{k,\mu}x^{k+\mu}D^\mu a\) with \(c_{k,0}=1\),
\(c_{k,1}=k\). For \(k<0\) one starts from \(ax^{-1}=x^{-1}a-Da\), which
is \([t,a]=Da\). Each term \(x^\mu D^\mu\) raises the power of \(x\) by
\(\mu\) and lowers the order by \(\mu\), so \(J+p\) is additive.
\end{proof}

Use the mixed complex of Lemma~\ref{ha:lem:mixed}, with \(\mathfrak
B_\partial=\C^Z[[s]][[x]]\) and stages \(S_{L,M}\) (coefficients in
\(x^{-L}s^{-M}\C^Z[[s]][[x]]\)). A chain has a \emph{bidegree}
\((J,p)\): the total power of \(x\) and the total order. By (F3),
\(S_{L,M,k}\) is the product of its bihomogeneous components, and each
component space is finite dimensional (the exponents in slot \(0\) are
\(\ge-L\), \(\ge-M\), the others \(\ge0\)). By
Lemma~\ref{ha:lem:n1alg}, \(b=\sum_{\mu\ge0}b_\mu\) with \(b_\mu\) of
bidegree \((\mu,-\mu)\). So \(S_{L,M}=\prod_gS^g_{L,M}\), where
\(S^g\) consists of the chains with \(J+p=g\), and on \(S^g\) the power of
\(x\) determines the bidegree. Filter \(S^g_{L,M}\) by \(J\). Each
\(S^g_{L,M}\) is the product of its finite-dimensional pieces of fixed
\(J\), with \(J\) bounded below, so this filtration is complete.

For forms, give \(x^as^r\,dx^{\epsilon_1}ds^{\epsilon_2}\) the bidegree
\((a+\epsilon_1,-r-\epsilon_2)\). Let \(\Omega(J,p)\) be the span of the
monomial forms of bidegree \((J,p)\) (on each \(z\)), and
\(\Omega_{L,M}(J,p)\) the span of those with \(a\ge-L\), \(r\ge-M\).
Call \((J,p)\) \emph{full} for \((L,M)\) if \(J\ge-L+1\) and \(p\le M-1\).
Then \(\Omega_{L,M}(J,p)=\Omega(J,p)\) for full \((J,p)\).

\begin{lemma}\label{ha:lem:n1E1}
\textup{(a)} \(\chi\) identifies the \(b_0\)-homology of \(S_{L,M}\) in
bidegree \((J,p)\) with \(\Omega^k_{L,M}(J,p)\), and \(b_1\) induces
\(-i\delta\), where \(\delta\) is the Koszul--Brylinski differential of
\eqref{eq:poisson} (\(n=1\)).

\textup{(b)} \(\delta\) has bidegree \((1,-1)\). If \(\theta\in\Omega^k(J,p)\)
satisfies \(\delta\theta=0\) and \((J,p)\ne(1-k,k-1)\), then
\(\theta=\delta\zeta\) with \(\zeta\in\Omega^{k+1}(J-1,p+1)\). In bidegree
\((1-k,k-1)\) the \(\delta\)-homology is \(H^{2-k}(Z\times\Ssig\times\Sd)\)
(on each \(z\), the images under \(*^{-1}\) of \(1\), \(dx/x\),
\(ds/s\), \(dx/x\wedge ds/s\)), as in Proposition~\ref{prop:Ads}.
\end{lemma}

\begin{proof}
(a) The \(x\)-graded algebra of \(\mathfrak B_\partial\) is the commutative
algebra \(R=\C^Z[[s]][x]\), and the graded coefficient module is \(uR\),
\(u=x^{-L}s^{-M}\), free of rank one with equal actions. So the graded
complex is the Hochschild complex of \(R\) shifted by \(u\), as in
Lemma~\ref{ha:lem:stageE1}. HKR holds for \(\C^Z[[s]]\)
(Theorem~\ref{ha:thm:variants}, formal series over the finite set \(Z\))
and survives tensoring with \(\C[x]\) (Proposition~\ref{ha:prop:kunneth}).
Since \(b_0\) and \(\chi\) preserve the bidegree, the identification holds
in each bidegree. For \(d^1\): the expansion of Lemma~\ref{ha:lem:n1alg}
has \(P_0(f,g)=fg\) and \(P_1(f,g)=x^2(\partial_xg)(Df)\) (with \(D\)
acting on the \(s\)-variable), so \(P_1(f,g)-P_1(g,f)=-i\{f,g\}\) with
\(\{f,g\}=\partial_\tau f\,x^2\partial_xg-x^2\partial_xf\,\partial_\tau g\).
Proposition~\ref{prop:d1} applies, as in the proof of
Theorem~\ref{ha:thm:conv}.

(b) The Poisson bivector \(x^2\partial_\tau\wedge\partial_x\) has
bidegree \((1,-1)\), and so does \(\delta\). The symplectic form
\(\omega=d\tau\wedge dx/x^2=\mp ds\wedge dx/(s^2x^2)\) has bidegree
\((-1,1)\). In dimension two \(*1=\omega\), \(*\omega=1\), and \(*\) maps
\(dx\) and \(ds\) to constant multiples of themselves (the computation in
the proof of Proposition~\ref{prop:Idpoint}). So \(*\) maps
\(\Omega^k(J,p)\) onto \(\Omega^{2-k}\bigl((J,p)+(k-1)(1,-1)\bigr)\). The de
Rham differential preserves the bidegree. On \(\Omega(J',p')\),
\(\mathcal L_{x\partial_x}=J'\) and \(\mathcal L_{s\partial_s}=-p'\), so by
Cartan's formula \(\iota_{x\partial_x}/J'\) (if \(J'\ne0\)) or
\(-\iota_{s\partial_s}/p'\) (if \(p'\ne0\)) is a contracting homotopy
that preserves the bidegree. For \((J',p')=(0,0)\) the forms \(1\),
\(dx/x\), \(ds/s\), \(dx/x\wedge ds/s\) are closed and not exact. Now
\(\delta=(-1)^{k-1}*^{-1}d*\) on \(k\)-forms (Proposition~\ref{prop:star}).
If \(\delta\theta=0\) and \((J,p)+(k-1)(1,-1)\ne(0,0)\), then
\(*\theta=d\eta\) with \(\eta\) of the same bidegree as \(*\theta\),
and \(\zeta=(-1)^k*^{-1}\eta\) is a \((k+1)\)-form with
\(\delta\zeta=(-1)^k*^{-1}d(*\zeta)=*^{-1}d\eta=\theta\). Its bidegree is that of \(\eta\) minus
\(k(1,-1)\), namely \((J-1,p+1)\).
\end{proof}

\begin{proposition}[Convergence for \(n=1\)]\label{ha:prop:xadic}
Let \(n=1\). Then \(\Theta:\HH_k(\Ads)\to E^2_{k-1,k}(\Ads)\) is an
isomorphism for every \(k\). Thus \(\HH_k(\Ads)\cong
H^{2-k}(Z\times\Ssig\times\Sd)\), and \(\HH_k(\Ads)=0\) for \(k\ge3\).
\end{proposition}

\begin{proof}
\emph{Step (P\('\)).} Fix \(L,M,g\) and a cycle \(c\in S^g_{L,M,k}\) whose
components have \(J\ge J_0\), where \((J_0-1,g-J_0+1)\) is full. Suppose
that \(g\ne0\) or \(J_0\ne1-k\). The component \(c_{J_0}\) is a
\(b_0\)-cycle, and \(b_1c_{J_0}+b_0c_{J_0+1}=0\). So
\(\theta=\chi(c_{J_0})\) satisfies \(\delta\theta=0\), and by
Lemma~\ref{ha:lem:n1E1}(b) \(\theta=-i\delta\zeta\) with \(\zeta\) of
bidegree \((J_0-1,g-J_0+1)\). This bidegree is full, so \(\zeta\in
\Omega_{L,M}\), and by Lemma~\ref{ha:lem:n1E1}(a) there is a
\(b_0\)-cycle \(x'\in S^g_{L,M}\) of power \(J_0-1\) with
\(\chi(x')=\zeta\). Then \(c_{J_0}-b_1x'\) has zero class in the stage
\(E^1\), so it equals \(b_0y'\) with \(y'\in S^g_{L,M}\) of power \(J_0\).
The cycle \(c-b(x'+y')\) has components with \(J\ge J_0+1\): the
components of \(b x'\) and \(by'\) of power \(J_0-1\) and \(J_0\) are
\(b_0x'=0\) and \(b_1x'+b_0y'=c_{J_0}\). The bidegree
\((J_0,g-J_0)\) is again full.

\emph{Off the critical part.} If \(g\ne0\), repeating (P\('\)) gives
corrections \(x'+y'\) of increasing power of \(x\). Their sum \(e\)
converges in the complete space \(S^g_{L,M,k+1}\), and \(c=be\), since
each component of \(be\) involves finitely many terms. If \(g=0\), the
same repetition brings \(c\) to a homologous cycle with components of
power \(J\ge1-k\), i.e.\ of order \(p=-J\le k-1=k-n\), as long as
\(J_0\le1-k\); if \(J_0>1-k\) it shows that \(c\) is a boundary.

\emph{Isomorphism.} Let \([c]\in\HH_k(\Ads)\). By
Lemma~\ref{ha:lem:mixed}(b), \(c\) can be taken in a stage \(S_{L,M}\);
we regard it as a chain of \(S_{L+2,M+2}\). All its components then lie
in bidegrees \((J,p)\) with \((J-1,p+1)\) full for \((L+2,M+2)\). The
components with \(g\ne0\) are boundaries in \(S_{L+2,M+2}\), and the sum
of their primitives over \(g\) is an element of the product
\(S_{L+2,M+2,k+1}\). So \([c]\) is represented by a cycle \(c'\in S^0\)
with components of order \(\le k-1\), whose leading term is the
component of bidegree \((1-k,k-1)\); this is
Proposition~\ref{ha:prop:laurent}(1), and \(\Theta[c]\) is the class of
that component. If \(\Theta[c]=0\), the class \(\theta\) of that
component is \(-i\delta\zeta\) with \(\zeta\in\Omega(-k,k)\). This
bidegree is full, because the component lies at power \(1-k\ge J_0\ge-L\)
and order \(k-1\le-J_0\le M\), where \(J_0\) is the lowest power of
\(c^0\) (if \(J_0>1-k\), then \(c^0\) is already a boundary). One more step (P\('\)) followed by the off-critical repetition
shows \(c'=be\). So \(\Theta\) is injective. For surjectivity, let
\(\theta\in\Omega^k(1-k,k-1)\) with \(\delta\theta=0\). Choose \(L,M\) such
that \((-k,k)\) and \((1-k,k-1)\) are full, and a \(b_0\)-cycle
\(x'\in S^0_{L,M}\) of power \(1-k\) with \(\chi(x')=\theta\). Since
\(\delta\theta=0\), \(b_1x'=-b_0y'\) with \(y'\) of power \(2-k\) (the
stage \(E^1\) is \(\Omega_{L,M}\) in every bidegree). The \((k-1)\)-cycle
\(b(x'+y')\) has components of power \(\ge3-k\), so it lies off the
critical power \(2-k\) of \((k-1)\)-chains, and the repetition above
writes it as \(be\) with \(e\) of power \(\ge2-k\). Then \(x'+y'-e\) is a
cycle with \(\Theta=[\theta]\).

The \(E^2\) term is that of Lemma~\ref{ha:lem:n1E1}(b), and it is the
\(E^2\) of Proposition~\ref{prop:Ads}, because the two computations use
the same forms and the same \(\delta\).
\end{proof}

\begin{corollary}[(B10) for \(n=1\)]\label{ha:cor:n1}
For \(n=1\), \(\Theta\) is an isomorphism for \(\As\) and for \(\Ads\).
Propositions~\ref{prop:As}, \ref{prop:Ads} and~\ref{prop:boundaryLES}
hold for \(n=1\) without quoted input.
\end{corollary}

\begin{proof}
Proposition~\ref{ha:prop:xadic} and Proposition~\ref{ha:prop:equiv}.
\end{proof}

\begin{remark}[The \(x\)-adic filtration for \(n\ge2\)]\label{ha:rem:xadicfail}
For \(n\ge2\) the \(x\)-adic filtration of \(\Ads\) still exists, and it
is the starting point of the proof in Appendix~\ref{su:sec}. Each
correction in the product raises the power of \(x\) by some \(\mu\ge0\)
and lowers the order by at least \(\mu\): the twist contributes
\((\mu,-\mu)\) as in Lemma~\ref{ha:lem:n1alg}, and the star product of
the suspended symbols \((0,-l)\). So the sum of the power of \(x\) and the
order never increases, and as the power of \(x\) grows the order is forced
down. The \(x\)-graded algebra is \(\mathcal S\otimes\C[x,x^{-1}]\), with
\(\mathcal S=\Psis^{\Z}(\partial X)/\Psis^{-\infty}(\partial X)\) the
suspended symbol algebra and \(x\) central. It is not commutative for
\(n\ge2\). Running the proof of Proposition~\ref{ha:prop:xadic} with the
\(x\)-adic filtration alone would need its \(E^1\) term
\(\HH(\mathcal S)\otimes\Omega(\C[x,x^{-1}])\) and the action of the twist
on it. That is more than is needed, and more than is easy: the order
spectral sequence of \(\mathcal S\) has \(E^1\) the homogeneous forms on
\((T^*\partial X\times\R_\tau)\setminus0\) and \(d^1=-i\delta\) for the
Poisson structure of \(T^*\partial X\), of which \(\tau\) is a Casimir;
that structure is not symplectic, and its \(E^2\) term is not
concentrated on one line (Lemma~\ref{su:lem:E2S},
Remark~\ref{su:rem:HS}). In the order filtration of \(\Ads\), on the
other hand, the primitives of Lemma~\ref{ha:lem:WL} lower the lowest power
of \(x\), and the terms \((0,-l)\) carry that power down to the next
order, so the towers of Proposition~\ref{ha:prop:laurent}(3) can have
unbounded poles. Appendix~\ref{su:sec} combines the two filtrations, as
Lauter and Moroianu do \cite{LauterMoroianu}: it runs the order
filtration, measures poles by the total power of \(x\) (so that a step
loses at most one pole), and removes the part with the extra pole using
the \(x\)-graded algebra. For that it needs only that the order spectral
sequence of \(\mathcal S\) degenerates at \(E^2\)
(Proposition~\ref{su:prop:degS}), not \(\HH(\mathcal S)\).
\end{remark}

\subsection{The \texorpdfstring{\(x\)}{x}-adic spectral sequence of
\texorpdfstring{\(\Id\)}{Id}}\label{ha:subsec:Id}

The proof of Proposition~\ref{prop:Idpoint} assumed that the corrections
of the \(x\)-adic filtration of \(\Id\) can be kept in a fixed pole stage.
The margin argument of Proposition~\ref{ha:prop:xadic} proves this, once
a complete model is available. Let \(\partial X\) be a point. Let
\(\mathcal P\) (the algebra \(\Ad\) of a boundary point) be the algebra of series \(\sum_{j\ge-L}x^ja_j\),
\(a_j\) classical symbols in \(\tau\) of order \(\le M\), with the
product of Lemma~\ref{ha:lem:n1alg} and \(D=-i\partial_\tau\)
(Lemma~\ref{lem:twist}). It is unital, contains
\(\Id^{\mathrm{pt}}=\cS(\R_\tau)\Lser{x}\) as an ideal, and has quotient
the algebra \(\Ads\) of Lemma~\ref{ha:lem:n1alg} (with \(Z\) two points).

\begin{proposition}[Convergence for \(\Id\)]\label{ha:prop:Idconv}
There is a complex \(D_\partial\) quasi-isomorphic to \(C(\Id^{\mathrm
{pt}})\) whose stages are complete for the \(x\)-adic filtration, whose
\(E^1\) is the space of forms on \(\R_\tau\times\Sd\), Schwartz in
\(\tau\) and formal Laurent in \(x\), and whose \(d^1\) is \(-i\delta\)
for \(\{x,\tau\}=-x^2\). The \(E^2\) term vanishes except at
\((J,k)=(1,0)\) and \((0,1)\), and taking the component of power
\(J=1-k\) of a suitable representative is an isomorphism
\(\HH_k(\Id^{\mathrm{pt}})\to E^2_{J=1-k,k}\).
\end{proposition}

\begin{proof}
\emph{The model.} Let \(\mathfrak B_d\subset\mathcal P\) be the
pole-free elements of order \(\le0\), a subalgebra by the product formula,
and \(\mathcal P_{L,M}=x^{-L}\{\text{order}\le M\}[[x]]\). With
\(T=\langle\tau\rangle^{-1}\in\mathfrak B_d\), which is invertible in
\(\mathcal P\) (elements of power \(0\) multiply pointwise), and
\(u=x^{-L}T^{-M}\), one has \(\mathcal P_{L,M}=\mathfrak B_du=u\mathfrak
B_d\), and \(\mathcal P_{N,N}\) has a continuous complement in
\(\mathcal P_{N+1,N+1}\). These are the only inputs of the proof of
Lemma~\ref{ha:lem:mixed}(b), which therefore applies: the mixed complex
\(C(\mathfrak B_d,\mathcal P)\) computes \(\HH(\mathcal P)\).
The full-symbol map to \(C(\mathfrak B_\partial,\Ads)\) is onto on every
stage (Borel's lemma in \(s\) at \(\tau=\pm\infty\)). Let \(D_\partial\) be
its kernel. \(\Id^{\mathrm{pt}}\) is H-unital (the contraction at the end
of the proof of Proposition~\ref{prop:Idpoint}), so, exactly as in the
proof of Theorem~\ref{ha:thm:conv}, excision and the five lemma show that
\(D_\partial\) computes \(\HH(\Id^{\mathrm{pt}})\).

\emph{Completeness and \(E^1\).} In a stage, the chains of fixed total
power \(J\) involve finitely many powers in the individual slots, and the
stage is the product over these, by (F3). So the \(x\)-adic filtration of
each stage of \(D_\partial\) (a closed subspace) is complete. Classical
symbols of order \(\le0\) are the smooth functions on the radial
compactification \(\overline\R\), a closed interval. By
Theorem~\ref{ha:thm:connes}, Proposition~\ref{ha:prop:kunneth} and the
shift by \(u\) (as in Lemma~\ref{ha:lem:n1E1}(a)), the stage \(E^1\) of
\(C(\mathfrak B_d,\mathcal P)\) in power \(J\) is \(u\) times the
forms of power \(J\) on \(\overline\R\times\) (formal \(x\)-line). The
full-symbol map induces Taylor expansion at the two ends of
\(\overline\R\), which is onto (Borel), so the stage \(E^1\) of
\(D_\partial\) is the kernel: forms that are flat at the ends, i.e.\
Schwartz in \(\tau\). The weight \(\langle\tau\rangle^M\) in \(u\) does
not change this space. In power \(J\ge-L+1\) all such forms of power
\(J\) occur. By Lemma~\ref{lem:twist} and Proposition~\ref{prop:d1},
\(d^1=-i\delta\), as in Lemma~\ref{ha:lem:n1E1}(a).

\emph{\(E^2\) with primitives one power lower.} As in
Lemma~\ref{ha:lem:n1E1}(b), \(*\) maps \(k\)-forms of power \(J\) to
\((2-k)\)-forms of power \(J+k-1\). The de Rham complex of forms Schwartz
in \(\tau\) and formal Laurent in \(x\) splits by the power of \(x\). In
power \(j\ne0\) it is contracted by the homotopy of
Lemma~\ref{lem:formaldR}, which preserves \(j\). In power \(0\) its
cohomology is \(H^\bullet_c(\R)\otimes\{1,dx/x\}\), in degrees \(1\) and
\(2\) (a Schwartz function with integral zero is the derivative of a
Schwartz function). Hence a \(\delta\)-closed \(k\)-form \(\theta\) of
power \(J\) with \((J,k)\ne(1,0),(0,1)\) is \(\delta\zeta\) with \(\zeta\)
of power \(J-1\), and the two exceptional groups are the ones computed in
Proposition~\ref{prop:Idpoint}.

\emph{Convergence.} Now repeat the proof of
Proposition~\ref{ha:prop:xadic} with the power of \(x\) alone (there is no
grading by the order, and none is needed): start from a cycle in a stage
\(S_{L,M}\), regard it in the stage with \(L+2\), and perform the steps
(P\('\)). Every primitive has power at least \(-L-1\ge-(L+2)+1\), so it
lies in the stage, every remainder has a strictly higher lowest power,
and the series of corrections converges by completeness. The critical
powers are \(J=1-k\).
\end{proof}

\begin{remark}[Benameur--Nistor]\label{ha:rem:BN}
Benameur and Nistor \cite{BN} compute the Hochschild homology of algebras
\(\mathcal A_L(G)=\mathcal O(M)\bigl(\Psi^\infty(G)/\Psi^{-\infty}(G)
\bigr)\) of Laurent complete symbols on a groupoid \(G\) with units \(M\),
under their assumptions (13) and (14), and state that the spectral
sequence degenerates at \(E^2\) and converges \cite[Thm.~6, Lemma~8]{BN}.
For the cusp groupoid of \(X\) (base \(B\) a point) the assumptions hold:
(13) is automatic, and (14) holds because, once \(x^{-1}\) is adjoined,
the \(\mathcal O(X)\)-span of \(x^2\partial_x,\partial_y\) is that of
\(\partial_x,\partial_y\). We cannot use this result to close (B10), for
reasons (2) and (3) below; (1) is a caveat about \cite{BN}. (The gap is
now closed by a different route, Theorem~\ref{su:thm:B10}.)
\begin{enumerate}
\item \emph{Theorem 4 of \cite{BN}.} By definition
\(\mathcal A_L(G)=\mathcal O(M)\bigl(\Psi^\infty(G)/\Psi^{-\infty}(G)\bigr)\);
for the cusp groupoid this is \(\As\). The b-groupoid of \(X\) also
satisfies (13)--(14), so \cite[Thm.~4]{BN} asserts that \(\As\) and its
b-analogue have the same image in
\(\Psi^\infty_{\mathrm{prop}}(X^\circ)/\Psi^{-\infty}_{\mathrm{prop}}(X^\circ)\).
This is false. The b-symbol \(|(\xi_b,\eta)|\) is
\(x^{-1}(\tau^2+x^2|\eta|^2)^{1/2}\) in cusp variables (\(\tau=x\xi_b\)),
and no \(x^L\) times it is smooth on \(\cT\setminus0\) at \(x=\tau=0\), so
it is not the principal symbol of an element of \(\As\). The gap is the
last step of the proof of \cite[Thm.~4]{BN}, which passes from equal
closures to equal ranges. For \(B\) a point, Theorem~4 enters \cite{BN}
only through the notation \(\mathcal A_L(F)\) in \cite[Lemma~7]{BN}, so this
error does not by itself affect their result for \(\As\); it means that
the input quoted from their earlier paper must be checked to concern the
cusp groupoid.
\item \emph{The base case is quoted.} For \(B\) a point, the case at
hand, \cite{BN} do not give a proof but refer to an earlier paper for
the \(E^2\) term and for \(\HH(\mathcal A_L(F))\) \cite[Lemma~5,
Lemma~7]{BN}.
\item \emph{Completeness.} The chains of \cite[(8)]{BN} are
\(\varinjlim_n(F^n_{n,n})^{\widehat\otimes(q+1)}\), stages in which the
pole bound and the order bound grow together, as in
Section~\ref{subsec:chains}. The proof of \cite[Thm.~1]{BN} uses
\(H(A)=\varprojlim_pH(A)/F_pH(A)\) for this inductive limit, and the
inductive limit is not complete (the example after
Proposition~\ref{ha:prop:algconv}). The complete subcomplexes that we
know, the stages of the mixed complex, have the stray \(E^2\) classes of
Example~\ref{ha:ex:obstruction}, so the vanishing hypothesis of
\cite[Thm.~1]{BN} cannot be checked on them.
\end{enumerate}
So \cite{BN} does not give (B10); they assert the corresponding statement
for \(\mathcal A_L(G)\), which for the cusp groupoid is \(\As\), as
\cite[p.~12]{MN} does. Lauter and Moroianu make the same point about
completeness. They note \cite[\S5]{LauterMoroianu} that, as observed by
Benameur and Nistor \cite{BNcomplete}, the spectral sequence argument
computes the \(E^2\) term and that it degenerates there, and that
convergence is the non-trivial part, because the pole filtration enters
the definition of the chain spaces. They prove convergence for the
fibered cusp algebras with their criterion \cite[Thm.~B.4]{LauterMoroianu}.
Appendix~\ref{su:sec} adapts that proof; it is how (B10) is proved in
these notes.
\end{remark}
